\newcommand{\DoNotLoadEpstopdf}{}
\documentclass[11pt]{article}
\usepackage[T1]{fontenc}
\usepackage{lmodern}
\usepackage[a4paper,margin=29mm]{geometry}
\usepackage{amsmath,amssymb,amsthm,mathtools}
\usepackage{booktabs,array}
\PassOptionsToPackage{hyphens}{url}
\usepackage[hidelinks]{hyperref}
\hypersetup{unicode=true,
  pdftitle={A negatively pinched K\"ahler threefold without bounded holomorphic coordinates},
  pdfauthor={OpenAI}}
\usepackage{microtype}
\ifdefined\pdfinfoomitdate\pdfinfoomitdate=1\fi
\ifdefined\pdftrailerid\pdftrailerid{}\fi
\ifdefined\pdfsuppressptexinfo\pdfsuppressptexinfo=15\fi
\newtheorem{theorem}{Theorem}[section]
\newtheorem{proposition}[theorem]{Proposition}
\newtheorem{lemma}[theorem]{Lemma}
\newtheorem{corollary}[theorem]{Corollary}
\theoremstyle{definition}
\newtheorem{definition}[theorem]{Definition}
\theoremstyle{remark}
\newtheorem{remark}[theorem]{Remark}
\newcommand{\B}{\mathbb B}
\newcommand{\Sph}{S}
\newcommand{\C}{\mathbb C}
\newcommand{\R}{\mathbb R}
\newcommand{\dd}{\mathrm d}
\DeclareMathOperator{\supp}{supp}
\DeclareMathOperator{\dist}{dist}
\DeclareMathOperator{\osc}{osc}
\numberwithin{equation}{section}
\allowdisplaybreaks[1]
\setlength{\emergencystretch}{1em}
\title{A negatively pinched K\"ahler threefold\\without bounded holomorphic coordinates}
\author{OpenAI}
\date{September 25, 2026}
\begin{document}
\maketitle
\begin{abstract}
We construct a contractible domain in complex dimension three with a
complete K\"ahler metric whose real sectional curvatures lie between two
finite negative constants. It admits no bounded holomorphic map to
$\mathbb C^3$ with nowhere-vanishing Jacobian and is therefore not
biholomorphic to a bounded domain. This gives a negative answer to the
negatively pinched K\"ahler uniformization question.
\end{abstract}
\setcounter{tocdepth}{1}
\tableofcontents
\clearpage
\section{Introduction}
\label{sec:introduction}

A bounded domain in $\C^n$ has bounded holomorphic coordinate functions
whose differentials are independent everywhere. The negatively curved
uniformization problem asks whether curvature and completeness can
force a complex manifold to possess such coordinates. In complex
dimension one, a complete simply connected surface whose curvature is
bounded above by a negative constant is biholomorphic to the disk. In higher
dimension, the existence of bounded holomorphic coordinates is a
separate global analytic requirement.

An early formulation appears in H.~Wu's work on normal families
\cite[p.~195, question~(1)]{WuNormal1967}. He asks whether a complete
simply connected K\"ahler manifold with nonpositive real sectional
curvature and holomorphic sectional curvature bounded above by a
negative constant is biholomorphic to a bounded domain. Wu and Yau
state the two-sided real-sectional version with a bounded pseudoconvex
domain as conclusion \cite[Conjecture~4.3, p.~18]{WuYauSurvey2019}:
if a simply connected complete K\"ahler manifold has every real
sectional curvature between two negative constants, must it be
biholomorphic to a bounded domain in its complex dimension?

The unit ball is already too restrictive a candidate in higher
dimension. Mostow and Siu constructed compact negatively curved
K\"ahler surfaces not covered by the ball \cite{MostowSiu1980}.
Deraux reviews their construction and obtains three-dimensional
examples \cite{Deraux2005}.
Their lifted metrics on the universal covers have two-sided negative
sectional bounds by compactness. These non-ball results, however, do
not exclude a biholomorphism to some other bounded domain; that is
the stronger obstruction sought here.

Here and throughout, sectional curvature refers to the underlying real
Riemannian metric, and completeness means geodesic completeness. A
bounded holomorphic map is a holomorphic map with bounded image in its
Euclidean target.

\begin{theorem}\label{main:theorem}
There exist a contractible domain $M\subset\C^3$, a smooth complete
K\"ahler metric $g$ on $M$, and finite constants $0<A\le B$ such that
\[
 -B\le K_g(\sigma)\le -A<0
\]
for every real two-plane $\sigma\subset T_pM$ and every $p\in M$.
No bounded holomorphic map $F:M\to\C^3$ has a nowhere-vanishing Jacobian
determinant. In particular, $M$ is not biholomorphic to a bounded domain
in $\C^3$.
\end{theorem}

Theorem~\ref{main:theorem} gives a negative answer to the stated
uniformization question, and its curvature hypotheses also place the
example within Wu's earlier question. The example has two bounded
nonconstant holomorphic functions, supplied by its base coordinates.
Its analytic obstruction concerns a fixed triple of bounded functions
with independent differentials at every point.
The domain is nevertheless Stein: the construction supplies a smooth
strictly plurisubharmonic exhaustion (Lemma~\ref{metric:complete}), so Grauert's solution
of the Levi problem applies \cite[Theorem~2 and Section~3.4]{GrauertLevi1958}.
Thus the obstruction is to bounded holomorphic coordinates, not to
holomorphic convexity.

\subsection{Intrinsic metrics and bounded functions}

Two-sided negative real sectional curvature has strong consequences for
intrinsic complex geometry. Greene and Wu established the lower
comparison for the Bergman metric \cite{GreeneWu1979}; Wu and Yau recall that result and
prove the complementary upper bound
\cite[Theorems~4 and~6]{WuYauInvariant2020}. Under completeness and
simple connectivity, Wu and Yau then compare the Bergman,
Kobayashi--Royden, complete negative K\"ahler--Einstein, and original
metrics \cite[Corollary~7]{WuYauInvariant2020}. The Bergman metric
uses square-integrable holomorphic top-degree forms. These comparisons
do not supply a fixed tuple of bounded holomorphic functions with
independent differentials everywhere. Applied to the example in
Theorem~\ref{main:theorem}, they show that all these intrinsic metrics
can be complete and mutually comparable while no such tuple exists.

The existence of even one nonconstant bounded holomorphic function is
a related question with a different conclusion. The companion paper
\cite[Theorem~1.1]{OpenAIOneSided2026} constructs, in a sufficiently
large fixed finite complex dimension, a domain diffeomorphic to
Euclidean space with a complete K\"ahler metric of real sectional
curvature at most $-1$ and only constant bounded holomorphic functions.
Its curvature is unbounded below. Here both curvature bounds are
finite and uniform, while the two bounded base coordinates remain.
The two constructions and their proofs are independent.

\subsection{Construction and proof strategy}

We construct $M$ as a family of discs over the unit ball
$\B=\{z\in\C^2:|z|<1\}$. A smooth real function $\phi$ on $\B$ determines
\[
 M_\phi=\{(z,w)\in\B\times\C:e^{\phi(z)}|w|^2<1\}.
\]
Thus the disc above $z$ has radius $e^{-\phi(z)/2}$. Put
$\psi(z)=-\log(1-|z|^2)$. Our metric is the complex Hessian of
\[
 \lambda\psi(z)-\log\bigl(1-e^{\phi(z)}|w|^2\bigr),
\]
with the real-metric convention specified in
Section~\ref{sec:reduction}. This is a specialization of Calabi's
construction from a base potential and a radial function of a Hermitian
fiber norm \cite[Section~3, equations~(3.1)--(3.2)]{Calabi1979}.
The momentum formulation of Hwang and Singer develops this framework
for general base and bundle data \cite{HwangSinger2002}.

For the logarithmic radial potential used here, Hao, Chen, and Wang
prove holomorphic and real sectional pinching results when the bundle
curvature determines the base metric
\cite[Theorems~1 and~2]{HaoChenWang2025}. Hao, Chen, Wang, and Zhang
obtain related results for complete K\"ahler--Einstein disk-bundle
metrics under Einstein and curvature hypotheses on the base
\cite[Sections~4.2--4.3]{HaoChenWangZhang2026}.
The estimate needed here concerns all real two-planes when the base
form $\lambda\partial\bar\partial\psi$ and the fiber-weight curvature
$\partial\bar\partial\phi$ are independent. We assume upper and lower
bounds for the latter form and bounds for the component derivatives
needed for curvature, expressed in centered ball coordinates. No curvature sign for the metric
$\partial\bar\partial\phi$ is assumed. We derive and estimate the
full curvature tensor under precisely these hypotheses.

The analytic task is to choose $\phi$ so that it has these derivative
bounds but cannot dominate a holomorphic real part, even after adding
the logarithmic boundary allowance forced by Cauchy's estimate for a
Jacobian. To do this, we construct radii $r_N\uparrow1$ and positive
probabilities $\mu_N$ on the unit sphere $S\subset\C^2$ such that
\[
 \int_S\left(\phi(r_N\xi)+4\log\frac1{1-r_N}\right)
                  \,d\mu_N(\xi)\longrightarrow-\infty,
\]
while $\int_S\operatorname{Re}H(r_N\xi)\,d\mu_N(\xi)
=\operatorname{Re}H(0)$ for every holomorphic $H$ on $\B$.
The two averages exclude any holomorphic real part bounded above by
the tested function plus a constant. The probabilities preserve
holomorphic evaluation at the origin, but such
measures need not preserve the logarithmic mean inequality for
holomorphic functions: Hedenmalm gives an explicit example of this
distinction \cite[Theorem~3.7]{Hedenmalm1989}.

The construction proceeds at degrees $k_j=Q^j$ for one fixed large
integer $Q$. At each degree we sum homogeneous peaks centered on a
separated set of complex lines. The projective packing and power-sum
estimate have a direct antecedent in Ryll and Wojtaszczyk
\cite[Section~2]{RyllWojtaszczyk1983}. A radial correction around each
peak annihilates every holomorphic moment exactly. Separation of the
degrees makes the corrections compatible with positivity at every
stage. Regularized logarithms built from the peak polynomials then have a
fixed negative average, which survives all subsequent corrections.
Summing these logarithms with a sufficiently large fixed coefficient
makes their negative averages dominate the boundary allowance.
Low degrees vary little
in a centered chart, while high degrees are exponentially small;
this gives the uniform derivatives through order four needed for
curvature.

The geometric task has its own delicate point. To prove negative
curvature, we must show that its negated numerator is positive on every
real two-plane. After rescaling the base directions, the leading
expression is nonnegative but can vanish on planes with nonzero base
and fiber projections. On these limiting planes the contribution of the
potentially largest error term vanishes as well. This cancellation gives stronger error
control near the degenerate planes and proves strict negativity for
one sufficiently large base parameter $\lambda$.
With $\lambda$ fixed, compactness on bounded
fiber ranges and a uniform limiting tensor at the fiber boundary give
both finite curvature bounds.

Section~\ref{sec:reduction} states the two precise ingredients and derives
Theorem~\ref{main:theorem} from them by Cauchy estimates for a putative
Jacobian. Sections~\ref{sec:peaks}--\ref{sec:densities} construct the
peaks and positive representing densities;
Sections~\ref{sec:tests}--\ref{sec:base-potential} establish their
logarithmic tests and the controlled potential.
Section~\ref{sec:hartogs} proves completeness.
Sections~\ref{sec:curvature}--\ref{sec:pinching} compute the curvature
and prove the two-sided estimate, including the limiting planes just
described.

\section{Two inputs and the bounded-map obstruction}\label{sec:reduction}

The construction separates a problem about functions on the ball from a
curvature estimate.  The first input is a strictly plurisubharmonic
function with controlled derivatives and an obstruction to a prescribed
holomorphic real-part bound.  The second turns its complex Hessian bounds
into a complete metric with uniformly negative real sectional curvature.
We state both inputs precisely and then prove their implication for
Theorem~\ref{main:theorem}.

\subsection{Centered charts and the two propositions}

Let $\B=\{z\in\C^2:|z|<1\}$ and put
\[
 \psi(z)=-\log(1-|z|^2),\qquad
 I=(\psi_{i\bar j}).
\]
Here a complex Hessian is also viewed as a Hermitian form, linear in its
first argument.  To center coordinates at a point $z_0\in\B$, let
$r=|z_0|$, choose a unitary map $U$ with $U(r,0)=z_0$, and set
\begin{equation}\label{base:charts}
 T=U\circ T_r,\qquad
 T_r(\zeta)=\left(
 \frac{r+\zeta_1}{1+r\zeta_1},
 \frac{\sqrt{1-r^2}\,\zeta_2}{1+r\zeta_1}
 \right).
\end{equation}
These maps are ball automorphisms with $T(0)=z_0$.  The identity
\[
 1-|T_r(\zeta)|^2
 =\frac{(1-r^2)(1-|\zeta|^2)}{|1+r\zeta_1|^2}
\]
shows that $T^*I$ has the same expression in $\zeta$ as $I$ has in $z$:
the additional terms in $\psi\circ T$ are constant or pluriharmonic.
In particular, this matrix is the identity at $0$ and has zero first
derivatives there.  We call \eqref{base:charts} the centered ball charts.

For a smooth real function $\phi$ on $\B$, write
$h=\partial\bar\partial\phi$ and, in a centered chart $T$, write
\[
 h^T_{i\bar j}(\zeta)
 =\partial_{\zeta_i}\partial_{\bar\zeta_j}(\phi\circ T)(\zeta).
\]
The uniform control needed below is the existence of constants
$0<c\le C<\infty$ such that, for every centered chart,
\begin{equation}\label{base:control}
 c\,\mathrm{Id}\le h^T(0)\le C\,\mathrm{Id}
\end{equation}
and
\begin{equation}\label{base:jets}
 \left|\partial_k h^T_{i\bar j}(0)\right|
 +\left|\partial_{\bar\ell}h^T_{i\bar j}(0)\right|
 +\left|\partial_k\partial_{\bar\ell}h^T_{i\bar j}(0)\right|
 \le C
 \qquad(i,j,k,\ell\in\{1,2\}).
\end{equation}
The derivatives here are ordinary derivatives of the displayed component
functions.  In invariant form, \eqref{base:control} says $cI\le h\le CI$.

\begin{proposition}[A controlled potential]\label{base:existence}
There exist $\phi\in C^\infty(\B,\R)$ and a finite constant $C$ for
which \eqref{base:control} and \eqref{base:jets} hold with $c=1$, and
such that no $H\in\mathcal O(\B)$ and $C_H\in\R$ satisfy
\begin{equation}\label{base:no-minorant}
 \operatorname{Re}H(z)
 \le C_H+4\log\frac1{1-|z|}+\phi(z)
 \qquad(z\in\B).
\end{equation}
\end{proposition}

For any smooth real $\phi$ on $\B$, define its Hartogs domain by
\begin{equation}\label{metric:domain}
 M_\phi=\{(z,w)\in\B\times\C:t(z,w)<1\},
 \qquad t(z,w)=e^{\phi(z)}|w|^2.
\end{equation}
For $\lambda>0$ define the potential and its complex Hessian by
\begin{equation}\label{metric:potential}
 \mathcal K_\lambda=\lambda\psi-\log(1-t),\qquad
 \widehat g_\lambda=\partial\bar\partial\mathcal K_\lambda.
\end{equation}
Our real metric convention is
$g_\lambda=2\operatorname{Re}\widehat g_\lambda$: if a real vector $X$
has $(1,0)$ part $u$, then
$g_\lambda(X,X)=2\widehat g_\lambda(u,u)$.

\begin{proposition}[Hartogs metric transfer]\label{metric:transfer}
Let $\phi\in C^\infty(\B,\R)$ satisfy \eqref{base:control} and
\eqref{base:jets} for fixed $0<c\le C<\infty$.  There is a finite
$\lambda_0=\lambda_0(c,C)>0$ such that, for every
$\lambda\ge\lambda_0$, the domain $M_\phi$ in
\eqref{metric:domain} is contractible and the tensor
$g_\lambda=2\operatorname{Re}\widehat g_\lambda$ defined by
\eqref{metric:potential} is a smooth geodesically complete K\"ahler
metric.  For this fixed $\lambda$ there are constants
$0<A_\lambda\le B_\lambda<\infty$ such that
\[
 -B_\lambda\le K_{g_\lambda}(\sigma)\le-A_\lambda<0
\]
for every point of $M_\phi$ and every real two-plane $\sigma$ in its
tangent space.
\end{proposition}

The proof of Proposition~\ref{base:existence} occupies the analytic
construction below.  The global geometry and curvature calculation in
Sections~\ref{sec:hartogs}--\ref{sec:pinching} prove
Proposition~\ref{metric:transfer}.  We first establish the short argument
that connects their conclusions to bounded holomorphic maps.

\subsection{The Jacobian on the zero section}

\begin{lemma}\label{map:obstruction}
Let $\phi\in C^\infty(\B,\R)$ have the nonexistence property in
Proposition~\ref{base:existence}: no holomorphic $H$ and real constant
$C_H$ satisfy \eqref{base:no-minorant}.  Then no bounded holomorphic
map $F:M_\phi\to\C^3$ has a nowhere-vanishing complex Jacobian
determinant.
\end{lemma}

\begin{proof}
Suppose that $F=(F_1,F_2,F_3)$ is such a map, and choose $B\ge1$ so
that $|F_\ell|\le B$ on $M_\phi$ for every component.  Since
$M_\phi$ is open in $\C^3$, its coordinates $(z_1,z_2,w)$ are global.
The function
\[
 d(z)=\det\frac{\partial(F_1,F_2,F_3)}{\partial(z_1,z_2,w)}(z,0)
\]
is holomorphic and nowhere zero on $\B$.

Fix $z\in\B$ and put $\delta=1-|z|$.  The restriction
$F_\ell(\,\cdot\,,0)$ is bounded and holomorphic on the entire base
ball.  For $j=1,2$, the disk
$\{z+\zeta e_j:|\zeta|<\delta\}$ is contained in that ball.  Cauchy's
estimate, applied on disks with radii tending to $\delta$, gives
\begin{equation}\label{map:base-cauchy}
 |\partial_{z_j}F_\ell(z,0)|\le B\delta^{-1}.
\end{equation}
For this fixed $z$, the fiber in $M_\phi$ is the disk of radius
$e^{-\phi(z)/2}$.  Cauchy's estimate on this disk gives
\begin{equation}\label{map:fiber-cauchy}
 |\partial_wF_\ell(z,0)|\le B e^{\phi(z)/2}.
\end{equation}
In particular, the base estimate requires no uniform fiber radius near
$z$; it is taken entirely on the zero section.

The six terms in the determinant expansion, together with
\eqref{map:base-cauchy} and \eqref{map:fiber-cauchy}, yield
\[
 |d(z)|\le6B^3\delta^{-2}e^{\phi(z)/2},
\]
and hence
\begin{equation}\label{map:determinant-bound}
 \log|d(z)|^2
 \le 2\log6+6\log B+4\log\frac1{1-|z|}+\phi(z).
\end{equation}
Because $\B$ is simply connected and $d$ has no zeros, it has a global
holomorphic logarithm $\ell$.  For example, a primitive of the closed
holomorphic one-form $d^{-1}\partial d$, with a suitable constant,
satisfies $e^\ell=d$.  The holomorphic function $H=2\ell$ has
$\operatorname{Re}H=\log|d|^2$, so \eqref{map:determinant-bound}
contradicts \eqref{base:no-minorant}.
\end{proof}

\begin{proof}[Proof of Theorem~\ref{main:theorem}]
Choose $\phi$ from Proposition~\ref{base:existence}, and then choose
$\lambda$ as in Proposition~\ref{metric:transfer}.  The domain
$M=M_\phi\subset\C^3$ is contractible and carries the asserted smooth
complete K\"ahler metric with two uniform negative real-sectional
bounds.  Lemma~\ref{map:obstruction} excludes a bounded holomorphic
map to $\C^3$ with nowhere-vanishing Jacobian.  A biholomorphism to a
bounded domain would be such a map: the chain rule applied to its
holomorphic inverse makes its Jacobian invertible everywhere.
\end{proof}

The logarithm argument uses nonvanishing along the entire zero section.
It gives no assertion that a bounded map's determinant vanishes
identically.  Nor are all bounded holomorphic functions constant:
$z_1$ and $z_2$ are nonconstant bounded holomorphic functions on $M_\phi$.

\section{Separated homogeneous peaks}\label{sec:peaks}

We begin the construction of the potential in
Proposition~\ref{base:existence}.  Its analytic requirement will follow
from increasingly negative averages against positive measures that
preserve holomorphic means.  The following observation explains the
role of those measures before we construct them.

Write $S=\{\xi\in\mathbb C^2:|\xi|=1\}$, and let $\sigma$ be its
surface measure normalized to have mass one.  A probability measure
$\mu$ on $S$ \emph{represents holomorphic evaluation at zero} if
\[
 \int_S p(\xi)\,d\mu(\xi)=p(0)
 \quad\text{for every holomorphic polynomial }p.
\]

\begin{lemma}[Representing-mean criterion]\label{peaks:mean-criterion}
Let $u$ be a continuous real function on $\mathbb B$.  Suppose there
are radii $0<r_N<1$ and representing probabilities $\mu_N$ on $S$
such that
\[
 \int_S u(r_N\xi)\,d\mu_N(\xi)\longrightarrow-\infty.
\]
Then no holomorphic function $H$ on $\mathbb B$ and constant $C_H$
satisfy $\operatorname{Re}H\le C_H+u$ on $\mathbb B$.
\end{lemma}

\begin{proof}
For each fixed $r_N<1$, the Taylor polynomials of $H$ converge
uniformly on $r_NS$.  The representing identities therefore give
\[
 \int_S\operatorname{Re}H(r_N\xi)\,d\mu_N(\xi)
 =\operatorname{Re}H(0).
\]
Integrating the proposed inequality contradicts the assumed limit.
\end{proof}

We will apply this criterion to
$u(z)=\phi(z)+4\log(1/(1-|z|))$.  The first ingredient is a supply
of homogeneous polynomials with many separated peaks.  Their uniform
bounds will also control the derivatives of the eventual potential.

\subsection{Projective packing and polynomial bounds}

The projective distance and separated-power estimate used here occur
in the direct construction of Ryll and Wojtaszczyk
\cite[p.~112 and Lemma~2.7]{RyllWojtaszczyk1983}.
We give the packing argument with a Gaussian tail and an adjustable
spacing parameter, since the density corrections will require both.
Related constructions with peaks localized in nonisotropic sphere
balls appear in Aleksandrov \cite[Lemmas~2--3]{Aleksandrov1983};
the caps used here are invariant under a common phase rotation.

Identify complex lines in $\mathbb C^2$ with points of
$\mathbb {CP}^1$.  For unit representatives define
\[
 d([\xi],[p])=\sqrt{1-|\langle\xi,p\rangle|^2},
\]
where the Hermitian inner product is linear in its first argument.
This equals the operator norm distance between the rank-one
orthogonal projections onto the two lines: their difference has
trace zero and eigenvalues
$\pm\sqrt{1-|\langle\xi,p\rangle|^2}$.  The triangle inequality
for the operator norm therefore proves that $d$ is a metric.  We also write
$d(\xi,p)$ when the representatives are understood.

In unitary coordinates with $\xi_1=\langle\xi,p\rangle$, write
\[
 \xi=(e^{i\theta}\cos v,e^{i\gamma}\sin v),
 \qquad 0\le v\le\pi/2.
\]
The round volume element is
$\cos v\sin v\,dv\,d\theta\,d\gamma$.  After normalization and
the change of variable $u=\cos^2v$, this gives
\begin{equation}\label{peaks:sphere-measure}
 d\sigma=du\,\frac{d\theta}{2\pi}\frac{d\gamma}{2\pi},
 \qquad 0\le u\le1.
\end{equation}
In particular, a projective cap $d(\xi,p)<s$ has measure $s^2$
for $0\le s\le1$.

Fix $D_0\ge1$.  For an integer $k\ge D_0^2$, choose a maximal
$D_0/\sqrt{k}$-separated set of lines, and choose one unit
representative $p$ for each line.  Let $\mathcal P_k$ denote this
set of representatives and $N_k=|\mathcal P_k|$.  Define the
homogeneous holomorphic polynomial
\begin{equation}\label{peaks:polynomial}
 P_k(z)=\sum_{p\in\mathcal P_k}\langle z,p\rangle^k.
\end{equation}
The representatives may be chosen arbitrarily; the estimates below
hold for every such choice.

\begin{lemma}[Peak bounds]\label{peaks:bounds}
There are absolute constants $C,c>0$, independent of $D_0\ge1$,
such that
\begin{equation}\label{peaks:count}
 \frac{k}{D_0^2}\le N_k\le\frac{4k}{D_0^2}
\end{equation}
and, for every $\xi\in S$ and $u\ge0$,
\begin{equation}\label{peaks:gaussian}
 \sum_{\substack{p\in\mathcal P_k\\\sqrt{k}d(\xi,p)\ge u}}
       |\langle\xi,p\rangle|^k\le Ce^{-cu^2}.
\end{equation}
Consequently,
\begin{equation}\label{peaks:global}
 |P_k(z)|\le C|z|^k\quad(z\in\mathbb C^2),
 \qquad
 \int_S|P_k|\,d\sigma\le\frac{2N_k}{k+2}.
\end{equation}
\end{lemma}

\begin{proof}
Put $\rho=D_0/\sqrt{k}\le1$.  Maximality gives a covering by caps
of radius $\rho$, and separation gives disjoint interiors for caps
of radius $\rho/2$.  Their measures yield
$1\le N_k\rho^2$ and $N_k\rho^2/4\le1$, proving
\eqref{peaks:count}.

We also need a local count whose constant does not depend on $D_0$.
About the centers within distance $u/\sqrt{k}$ of a fixed line,
place caps of radius $1/(3\sqrt{k})$.  Their interiors are disjoint,
and they lie inside the cap of radius $(u+1/3)/\sqrt{k}$ about that
line.  Comparing measures bounds their number by $9(u+1/3)^2$ when
the latter radius is at most one.  Otherwise their number is at
most $9k\le9(u+1/3)^2$, by comparison with the whole space.
Thus the local count is at most $C(1+u)^2$ in all cases.

Since
\[
 |\langle\xi,p\rangle|^k
 =(1-d(\xi,p)^2)^{k/2}
 \le e^{-kd(\xi,p)^2/2},
\]
the contribution from the annulus
$n\le\sqrt{k}d(\xi,p)<n+1$ is at most
$C(n+2)^2e^{-n^2/2}$.  Summing these bounds from $n=\lfloor u\rfloor$
onward proves \eqref{peaks:gaussian}: the polynomial factor is
absorbed by a weaker Gaussian exponent, and the range $0\le u\le2$
is absorbed by the constant.  Setting $u=0$ and using homogeneity
gives the first assertion in \eqref{peaks:global}, including points
outside the unit ball.  Finally, \eqref{peaks:sphere-measure} gives
\[
 \int_S|\langle\xi,p\rangle|^k\,d\sigma
 =\int_0^1 u^{k/2}\,du=\frac2{k+2}.
\]
The triangle inequality proves the stated integral bound.
\end{proof}

\subsection{Patches and a radial profile}

For a fixed $R_0>0$ and a center $p\in\mathcal P_k$, use the
adapted coordinates above and put
\begin{equation}\label{peaks:patch-definition}
 y=-k\log|\xi_1|,\qquad
 \mathcal U_{k,p}=\{\xi\in S:0\le y<R_0\}.
\end{equation}
The set $y=0$ is the phase circle through $p$.  Every patch is
invariant under multiplication of $\xi$ by a common phase.

\begin{lemma}[Patch estimates]\label{peaks:patches}
Fix $R_0>0$.  If $D_0>\max\{1,2\sqrt{2R_0}\}$ and
$k\ge\max\{D_0^2,2R_0\}$, the patches $\mathcal U_{k,p}$
are pairwise disjoint.  On each patch,
\begin{equation}\label{peaks:patch-measure}
 d\sigma=\frac2k e^{-2y/k}\,dy\,
             \frac{d\theta}{2\pi}\frac{d\gamma}{2\pi},
\end{equation}
and
\begin{equation}\label{peaks:single-peak}
 |P_k(\xi)-e^{-y}e^{ik\theta}|
 \le C e^{-c(D_0-\sqrt{2R_0})^2}.
\end{equation}
Outside their union,
\begin{equation}\label{peaks:off-patch}
 |P_k(\xi)|\le Ce^{-cR_0}.
\end{equation}
The constants $C,c$ can be chosen independently of $R_0,D_0,k$.
\end{lemma}

\begin{proof}
The projective radius of a patch is
$\sqrt{1-e^{-2R_0/k}}\le\sqrt{2R_0/k}$, so the separation of
the centers makes the patches disjoint.  Formula
\eqref{peaks:patch-measure} follows from
$|\xi_1|^2=e^{-2y/k}$ in \eqref{peaks:sphere-measure}.

If $\xi\in\mathcal U_{k,p}$ and $p'\ne p$ is another center,
the metric triangle inequality gives
\[
 \sqrt{k}d(\xi,p')\ge D_0-\sqrt{2R_0}.
\]
The peak from $p$ is exactly $e^{-y}e^{ik\theta}$.
Summing all other peaks by \eqref{peaks:gaussian} proves
\eqref{peaks:single-peak}.
Outside every patch, each center satisfies
\[
 kd(\xi,p)^2\ge k(1-e^{-2R_0/k})\ge R_0.
\]
The final inequality uses $k\ge2R_0$ and
$1-e^{-v}\ge v/2$ for $0\le v\le1$.  Another application of
\eqref{peaks:gaussian} proves \eqref{peaks:off-patch}.
\end{proof}

The density correction will use a radial profile with an exact
weighted cancellation.  The following choice also supplies the
strict inequality needed for the logarithmic tests later.

\begin{lemma}[Radial profiles]\label{peaks:profiles}
Fix $1<a<2$ and put $\alpha=\log a$.  One can choose $0<s_1<1$,
smooth real functions $f,b$ on $[0,\infty)$, and $R_0>0$ such that
\begin{equation}\label{peaks:profile-negativity}
 \alpha^2-2s_1(\cosh\alpha-1)<0,
\end{equation}
\begin{equation}\label{peaks:profile-identities}
 \|f\|_\infty<1,\quad f=-s_1\ \text{on }[0,\alpha],\quad
 \int_0^\infty e^{-y}f(y)\,dy=0,\quad
 b'-b=f,\quad b(0)=0.
\end{equation}
Both functions vanish on a neighborhood of $[R_0,\infty)$.
Moreover, $R_0$ can be chosen so that $|aP_k|\le1/2$ off the
patches whenever the hypotheses of Lemma~\ref{peaks:patches} hold.
\end{lemma}

\begin{proof}
The strict inequality $2(\cosh\alpha-1)>\alpha^2$ allows a choice
\[
 \frac{\alpha^2}{2(\cosh\alpha-1)}<s_1<1.
\]
Since $a<2$, this choice also satisfies
$s_1(1-e^{-\alpha})<e^{-\alpha}$.
Choose a small $\delta>0$ and a smooth cutoff $\chi$, equal to one
on $[0,\alpha]$, between zero and one, and zero for
$y\ge\alpha+\delta$.  For sufficiently small $\delta$,
\[
 N=s_1\int_0^\infty e^{-y}\chi(y)\,dy
 <e^{-(\alpha+2\delta)}.
\]
Choose a smooth compactly supported function $\nu$, supported
strictly after $\alpha+2\delta$, with $0\le\nu\le1$, and with
\[
 P=\int_0^\infty e^{-y}\nu(y)\,dy>N.
\]
Such a function is obtained by approximating the indicator of a
sufficiently long interval beginning just after $\alpha+2\delta$.
Then $f=-s_1\chi+(N/P)\nu$ has the required plateau and weighted
integral.  The two supports are disjoint, so
$\|f\|_\infty\le\max\{s_1,N/P\}<1$.

Define
\[
 b(y)=e^y\int_0^y e^{-v}f(v)\,dv.
\]
It satisfies $b'-b=f$ and $b(0)=0$.  Beyond the support of $f$,
the integral is exactly zero by the weighted cancellation, so $b$
also has compact support.  Near zero,
$b(y)=-s_1(e^y-1)$, which is smooth up to the boundary of the
half-line.  Smoothness here does not require extending $f$ by zero
to negative arguments.  Finally choose $R_0$ larger than both
supports, and enlarge it further until
$aCe^{-cR_0}\le1/2$ in \eqref{peaks:off-patch}.
\end{proof}

The plateau $0\le y\le\alpha$ is exactly the region where the
model peak $ae^{-y}$ has modulus at least one.  For the profiles
just constructed, the integral of
$\log|1-ae^{-y}e^{it'}|^2$ against
$(1+f(y)\cos t')\,dy\,dt'/(2\pi)$ on
$[0,R_0]\times[0,2\pi]$ equals
$\alpha^2-2s_1(\cosh\alpha-1)$, as calculated in
\eqref{tests:model-sign}.  Thus the strict profile inequality
prescribes a negative logarithmic integral, while the weighted
cancellation will preserve holomorphic moments.

Fix these profiles and $R_0$ from now on.  The spacing $D_0$ and
the degrees will remain available to be chosen sufficiently large.
In particular, enlarging $D_0$ later improves the single-peak
approximation without altering the profiles.

\section{Positive representing densities}\label{sec:densities}

We now construct the representing probabilities required by
Lemma~\ref{peaks:mean-criterion}.  Each step adds an oscillation on
the current patches.  A radial integration by parts makes every
holomorphic moment of that oscillation vanish exactly.  The
separation between consecutive degrees will ensure that the new
density stays positive.

Keep the profiles $f,b,R_0$ from Lemma~\ref{peaks:profiles} fixed,
and put
\[
 c_0=\frac{1-\|f\|_\infty}{4}>0.
\]
Fix any $D_0>\max\{1,2\sqrt{2R_0}\}$.  For an integer $Q\ge2$
large enough that every degree below meets the patch thresholds,
put
\begin{equation}\label{density:degrees}
 k_j=Q^j,\qquad \ell_0=0,\qquad
 \ell_j=\sum_{l=1}^j k_l\quad(j\ge1).
\end{equation}
Thus $\ell_{j-1}<k_j$ and $\ell_j<2k_j$.  At each degree use any
of the separated sets and patches constructed in Section~\ref{sec:peaks}.

Start with $W_0=1$.  At stage $j$, set $k=k_j$.  For each
$p\in\mathcal P_k$ let
\[
 S_p(\theta)=W_{j-1}(e^{i\theta}p).
\]
To determine the correction, consider a phase mode $e^{is\theta}$
with $|s|\le\ell_{j-1}$, the frequency range maintained below.
Multiplication by $e^{ik\theta}$ shifts it to frequency
$m=k+s>0$; its conjugate can pair with the holomorphic monomial
$\xi_1^m=e^{-my/k}e^{im\theta}$.  After canceling the surface-measure factor
$e^{-2y/k}$, its radial moment will vanish if its coefficient is
$b'-(m/k)b$, because
\[
 e^{-my/k}\left(b'-\frac{m}{k}b\right)
 =\frac{d}{dy}\left(e^{-my/k}b\right),
 \qquad b(0)=b(R_0)=0.
\]
The identity $f=b'-b$ therefore prescribes the coefficient
$f-(s/k)b$ for this mode.  Summing the modes leads to the following
definition in the adapted coordinates of $\mathcal U_{k,p}$:
\begin{equation}\label{density:recursion}
 \begin{split}
 A_{j,p}(\xi)
 &=e^{2y/k}\operatorname{Re}\left[
     e^{ik\theta}\left(f(y)S_p(\theta)
               -b(y)\frac{S_p'(\theta)}{ik}\right)\right],\\
 W_j&=W_{j-1}+\sum_{p\in\mathcal P_k}A_{j,p},
 \end{split}
\end{equation}
where each correction is extended by zero off its patch.  The
prime on $S_p$ denotes differentiation in $\theta$.  This recursion
is defined without assuming positivity; that property will be
proved for one fixed sufficiently large $Q$.

\begin{lemma}[Smoothness of the corrections]\label{density:smooth}
Every $W_j$ defined by \eqref{density:recursion} is a smooth real
function on $S$.
\end{lemma}

\begin{proof}
Proceed by induction, beginning with $W_0=1$.  The restriction
$S_p$ is a smooth periodic function.  To check the central circle
of a patch, use the smooth coordinates
\[
 (\theta,\zeta)\longmapsto
       (e^{i\theta}\sqrt{1-|\zeta|^2},\zeta),
 \qquad y=-\frac{k}{2}\log(1-|\zeta|^2).
\]
All terms of \eqref{density:recursion} are smooth in these
coordinates, including at $\zeta=0$.  They do not depend on the
transverse polar angle $\gamma$, which would be singular there.
Their dependence on $\theta$ is periodic, so it defines a smooth
function along the full central circle.  At the outer boundary,
$f$ and $b$ vanish in a collar of $R_0$; hence the correction is
identically zero in a collar of the patch boundary.  Its extension
by zero is therefore smooth.  Summing the finitely many corrections
proves the induction.
\end{proof}

Let $T$ denote the phase derivative on $S$,
\[
 TW(\xi)=\left.\frac{d}{d\omega}W(e^{i\omega}\xi)
                  \right|_{\omega=0}.
\]
A horizontal vector $Z\in T_\xi S$ will always mean a
\emph{Euclidean unit} real tangent vector satisfying the complex
equation $\langle Z,\xi\rangle=0$.

A function has phase bandwidth at most $\ell$ if, on every orbit
$\omega\mapsto e^{i\omega}\xi$, it is a trigonometric polynomial
with frequencies in $[-\ell,\ell]$.  Its orbitwise mean is the
constant coefficient of this polynomial.

\begin{theorem}[Uniform representing densities]\label{density:main}
There is a constant $K_*>0$, depending only on the fixed profiles
$f,b$ and $R_0$, with the following property.  For every fixed
$D_0>\max\{1,2\sqrt{2R_0}\}$ there is an integer threshold $Q_0$
such that, for every integer $Q\ge Q_0$ and every choice of the
separated sets above, the recursion \eqref{density:recursion}
has the following properties for all $j\ge1$:
\begin{enumerate}
\item $W_j$ is smooth and strictly positive, and $W_j\,d\sigma$
represents holomorphic evaluation at zero.  In particular it is
a probability measure.
\item Under global phase rotation, $W_j$ has frequencies only
between $-\ell_j$ and $\ell_j$, and its mean on every phase orbit
is exactly one.
\item At every point of $S$,
\begin{equation}\label{density:uniform-bounds}
 \begin{gathered}
 c_0W_{j-1}\le W_j\le2W_{j-1},\\
 |ZW_j|\le K_*\sqrt{k_j}\,W_j,
 \qquad |TW_j|\le K_*k_jW_j
 \end{gathered}
\end{equation}
for every unit horizontal vector $Z$.
\end{enumerate}
The threshold $Q_0$ may also be enlarged to impose any fixed lower
threshold on the degrees $k_j$.
\end{theorem}

We first prove the exact moment and frequency assertions; they
hold before positivity is known.  The remainder of the theorem,
including the order of choices of $K_*$ and $Q_0$, is proved in
the next subsection.

\subsection{Exact moments and phase frequencies}

\begin{lemma}[Moment preservation]\label{density:moments}
For every stage of \eqref{density:recursion}, the phase bandwidth
of $W_j$ is at most $\ell_j$, its orbitwise mean is one, and
\begin{equation}\label{density:representing}
 \int_S p(\xi)W_j(\xi)\,d\sigma(\xi)=p(0)
 \quad\text{for every holomorphic polynomial }p.
\end{equation}
Each individual correction has zero integral against every such
polynomial, including the constant polynomial.
\end{lemma}

\begin{proof}
At stage zero the phase assertions are immediate.  The polynomial
identity for $\sigma$ follows by global phase averaging: every
nonconstant homogeneous holomorphic polynomial has nonzero
positive phase frequency.

Assume the assertions through stage $j-1$, and write $k=k_j$.
The central restriction has a finite expansion
\[
 S_p(\theta)=\sum_{|s|\le\ell_{j-1}}\widehat S_p(s) e^{is\theta},
 \qquad \widehat S_p(0)=1.
\]
Since $b'-b=f$, the bracket in \eqref{density:recursion} is
\[
 \sum_s \widehat S_p(s)\left(b'-\frac{k+s}{k}b\right)e^{is\theta}.
\]
All integers $m=k+s$ are positive because $\ell_{j-1}<k$.
The patch and $y$ are invariant under global phase rotation,
whereas $\theta$ increases by the rotation angle.  Hence the
correction has frequencies in
\begin{equation}\label{density:bandwidth}
 [k-\ell_{j-1},k+\ell_{j-1}]
 \ \cup\ [-k-\ell_{j-1},-k+\ell_{j-1}].
\end{equation}
They omit zero.  Multiplying by a function of $y$ and extending
by zero across the invariant patch boundary do not alter these
frequencies.  The bandwidth of $W_j$ is thus at most
$k+\ell_{j-1}=\ell_j$, and its orbitwise mean remains one.

For the full moment calculation, use an adapted monomial
$\xi_1^n\xi_2^l$, with nonnegative integers $n,l$.  If $l>0$,
averaging in the transverse phase $\gamma$ gives zero.  If
$l=n=0$, averaging in $\theta$ gives zero because all $m=k+s$
are positive.  Suppose then that $l=0$ and $n\ge1$.  As $b$ is
real, the correction is
\[
 A_{j,p}=\frac{e^{2y/k}}2
 \sum_s\left(\widehat S_p(s) e^{i(k+s)\theta}
              +\overline{\widehat S_p(s)}e^{-i(k+s)\theta}\right)
       \left(b'-\frac{k+s}{k}b\right).
\]
Since $\xi_1^n=e^{-ny/k}e^{in\theta}$, only the conjugate
frequency with $k+s=n$ can survive the $\theta$ integration.
Formula~\eqref{peaks:patch-measure} gives exactly
\[
 \begin{split}
 \int_{\mathcal U_{k,p}}\xi_1^n A_{j,p}\,d\sigma
 &=\frac{\overline{\widehat S_p(n-k)}}{k}
   \int_0^{R_0}e^{-ny/k}\left(b'-\frac nk b\right)dy\\
 &=\frac{\overline{\widehat S_p(n-k)}}{k}
   \left[e^{-ny/k}b(y)\right]_0^{R_0}=0.
 \end{split}
\]
Here a coefficient outside the previous bandwidth is zero.  The
factor $2/k$ in surface measure combines with the $1/2$ from the
real part to leave $1/k$.  Both endpoint terms vanish because
$b(0)=b(R_0)=0$.
The argument applies to every degree $n$, without a degree cutoff.
Unitary changes of coordinates preserve holomorphic polynomials,
so each correction annihilates every polynomial in the original
coordinates.  Adding the corrections proves
\eqref{density:representing} and completes the induction.
\end{proof}

In particular, every $W_j$ has total integral one.  What remains
for Theorem~\ref{density:main} is to prove that these exact
representing densities are positive with the uniform derivative
bounds.  The next estimates accomplish this with a single choice
of $Q$, independent of the number of stages.

\subsection{Uniform positivity and derivative estimates}

We now prove the remaining assertions of Theorem~\ref{density:main}.
The main point is to choose one derivative constant \(K_*\), followed
by one threshold for \(Q\), that works at every stage.  A phase
derivative of the correction in \eqref{density:recursion} involves
\(S_p''\).  We obtain a bound for this derivative from the frequency
support and the estimates already proved at earlier stages.

\begin{lemma}[Growth along phase orbits]\label{density:orbit-growth}
Suppose \(Q\ge2\), \(k_j=Q^j\), and the positive functions
\(W_0=1,W_1,\ldots,W_n\), with \(n\ge1\), satisfy
\[
 c_0W_{j-1}\le W_j\le2W_{j-1},
 \qquad |TW_j|\le K_*k_jW_j
 \quad(1\le j\le n).
\]
Put
\[
 R=\frac2{c_0},\qquad \beta=\log_2R.
\]
For every \(p\in S\), the function
\(S_p(\theta)=W_n(e^{i\theta}p)\) satisfies
\begin{equation}\label{density:orbit-growth-bound}
 S_p(\theta+u)
 \le e^{K_*}R(1+k_n|u|)^\beta S_p(\theta),
 \qquad |u|\le\pi.
\end{equation}
In particular, the constants in this bound are independent of \(Q,n,p\).
\end{lemma}

\begin{proof}
Choose the largest \(1\le l\le n\) such that \(k_l|u|\le1\),
or put \(l=0\) if no such index exists.  The comparisons between
successive densities give
\[
 \frac{W_n(e^{i(\theta+u)}p)}{W_n(e^{i\theta}p)}
 \le R^{n-l}
 \frac{W_l(e^{i(\theta+u)}p)}{W_l(e^{i\theta}p)}.
\]
When \(l\ge1\), logarithmic differentiation and the phase derivative
bound show that the last ratio is at most
\(\exp(K_*k_l|u|)\le e^{K_*}\).
For \(l=0\) it equals one.
If \(l=n\), this already proves the assertion.
Otherwise the maximality of \(l\), including the case \(l=0\), gives
\[
 k_n|u|>Q^{n-l-1}\ge2^{n-l-1}.
\]
It follows that
\(R^{n-l}\le R(1+k_n|u|)^\beta\), proving
\eqref{density:orbit-growth-bound}.
\end{proof}

\begin{lemma}[A weighted derivative bound]\label{density:weighted-derivative}
Fix \(A\ge1\) and \(\beta\ge0\).  Let \(q\ge1\) be an integer,
and let \(s\) be a positive trigonometric polynomial of degree at most
\(2q\).  Suppose that
\[
 s(\theta+u)\le A(1+q|u|)^\beta s(\theta)
 \qquad(\theta\in\mathbb R,\ |u|\le\pi).
\]
There is a constant \(C(A,\beta)\), independent of \(q,s\), such that
\[
 |s''(\theta)|\le C(A,\beta)q^2s(\theta).
\]
\end{lemma}

\begin{proof}
Choose a fixed smooth compactly supported function
\(\chi:\mathbb R\to\mathbb R\) which equals one on \([-2,2]\), and put
\[
 K_q(u)=\sum_{h\in\mathbb Z}\chi(h/q)e^{ihu}.
\]
For normalized convolution on the circle, the frequency assumption
implies \(s''=s*K_q''\).  We claim that, for every positive integer \(N\),
\begin{equation}\label{density:kernel-decay}
 |K_q''(u)|\le C_Nq^3(1+q|u|)^{-N},
 \qquad |u|\le\pi,
\end{equation}
where \(C_N\) depends only on \(N,\chi\).

Indeed the coefficients of \(K_q''\) are
\[
 a_h=-h^2\chi(h/q)=q^2g(h/q),
 \qquad g(v)=-v^2\chi(v).
\]
There are \(O(q)\) nonzero coefficients, each of size \(O(q^2)\),
so direct summation gives \(|K_q''|\le Cq^3\).
An \(N\)-th finite difference of \(a_h\) is an iterated integral of
the \(N\)-th derivative of \(q^2g(v/q)\), and hence has size at most
\(C_Nq^{2-N}\).  Only \(O_N(q)\) such differences can be nonzero,
since \(g\) is compactly supported and \(q\ge1\).  Their absolute sum
is therefore at most \(C_Nq^{3-N}\).
Multiplication of the Fourier series by \((1-e^{iu})^N\) produces
these differences, up to an index shift.  Using
\(|1-e^{iu}|\ge2|u|/\pi\) on \([-\pi,\pi]\) gives the second estimate
\[
 |K_q''(u)|\le C_Nq^{3-N}|u|^{-N}\qquad(u\ne0).
\]
Combining it with the direct estimate proves
\eqref{density:kernel-decay}.

Choose an integer \(N>\beta+1\).
Positivity, the assumed growth bound, and
\eqref{density:kernel-decay} now give
\[
 \begin{split}
 |s''(\theta)|
 &\le\int_{-\pi}^{\pi}s(\theta-u)|K_q''(u)|
                         \,\frac{du}{2\pi}\\
 &\le AC_Nq^3s(\theta)
       \int_{-\pi}^{\pi}(1+q|u|)^{\beta-N}\,\frac{du}{2\pi}\\
 &\le C(A,\beta)q^2s(\theta).
 \end{split}
\]
\end{proof}

Combining these two lemmas with the bandwidth assertion of
Lemma~\ref{density:moments} gives the following consequence.
If the estimates of Theorem~\ref{density:main} hold through stage
\(j-1\), then for \(k=k_j\) and the central function \(S_p\) in
\eqref{density:recursion},
\begin{equation}\label{density:central-derivatives}
 \frac{|S_p'|}{kS_p}\le\frac{K_*}{Q},
 \qquad
 \frac{|S_p''|}{k^2S_p}\le\frac{C_1}{Q^2},
\end{equation}
where \(C_1\) depends only on \(K_*,c_0\) and the fixed cutoff
\(\chi\).  It is independent of \(Q,j,p\).
For \(j\ge2\), apply the lemmas with \(q=k_{j-1}\), noting that
\(\ell_{j-1}\le2k_{j-1}\); the first inequality follows directly
from the phase derivative bound.  For \(j=1\), both derivatives vanish
because \(S_p=1\).

We next estimate a single correction using only normalized bounds.
This will allow us to choose \(K_*\) before requiring
\eqref{density:central-derivatives} to be small.
All horizontal vectors \(Z\) below have Euclidean length one.
Enlarge a constant \(D_R\ge1\), depending only on \(R_0\), so that
on every patch at every sufficiently large degree \(k\),
\begin{equation}\label{density:patch-directions}
 |Yy|\le D_R,\qquad |Y\theta|\le D_R/k,\qquad
 |Y(e^{ik\theta})|\le D_R,
 \quad Y=Z/\sqrt{k}\ \hbox{or}\ Y=T/k.
\end{equation}
Here is a direct verification.  With \(c=\langle\xi,p\rangle\),
complex horizontality gives
\[
 |Zc|\le\sqrt{1-|c|^2}
 \le\sqrt{2R_0/k},\qquad |c|\ge e^{-R_0/k}.
\]
Differentiate \(y=-k\log|c|\) and \(\theta=\arg c\) to obtain
the first two bounds for \(Z/\sqrt{k}\); the third follows from
the second.  For \(T/k\), one has \(Yy=0\) and \(Y\theta=1/k\).
The fixed lower threshold for \(k\) makes \(e^{-R_0/k}\ge1/2\).
We also choose \(D_R\) large enough to bound the length of the radial
arc from \(e^{i\theta}p\) to each patch point by \(D_R/\sqrt{k}\).
In adapted coordinates this arc is
\[
 u\longmapsto(e^{i\theta}\cos u,e^{i\gamma}\sin u),
 \qquad 0\le u\le \arccos(e^{-y/k}).
\]
It has unit speed and is complex-horizontal.  The length bound follows
from \(0\le y\le R_0\) and \(1-e^{-y/k}\le R_0/k\).
Fix this enlarged \(D_R\) throughout the rest of the proof.

\begin{lemma}[Normalized correction estimate]
\label{density:normalized-patch}
Write
\[
 F=\|f\|_\infty,\quad B=\|b\|_\infty,\quad
 F_1=\|f'\|_\infty,\quad B_1=\|b'\|_\infty.
\]
On a degree-\(k\) patch suppose \(W_{j-1}>0\), \(S_p>0\), and
\[
 \frac{e^{2y/k}S_p}{W_{j-1}}\le2,\qquad
 r_1=\frac{|S_p'|}{kS_p}\le1,\qquad
 r_2=\frac{|S_p''|}{k^2S_p}\le1.
\]
For the correction \(A_{j,p}\) in \eqref{density:recursion},
\[
 |YA_{j,p}|\le C_2W_{j-1},
 \quad Y=Z/\sqrt{k}\ \hbox{or}\ T/k,
 \qquad
 C_2=2D_R\bigl[4(F+B)+F_1+B_1\bigr].
\]
In particular, \(C_2\) is independent of \(K_*,Q,j,p\).
\end{lemma}

\begin{proof}
Besides \eqref{density:patch-directions}, use
\[
 |Y(e^{2y/k})|\le(2D_R/k)e^{2y/k},\qquad
 |YS_p|\le(D_R/k)|S_p'|,\qquad
 |YS_p'|\le(D_R/k)|S_p''|.
\]
Differentiating \eqref{density:recursion} and taking absolute values gives
\[
 |YA_{j,p}|
 \le D_Re^{2y/k}S_p
 \left[(1+2/k)(F+Br_1)+F_1+B_1r_1+Fr_1+Br_2\right].
\]
Since \(k\ge1\), the assumed normalized bounds imply the stated estimate.
\end{proof}

\begin{proof}[Completion of the proof of Theorem~\ref{density:main}]
Fix \(K_*>2C_2/c_0\).  The constant \(C_1\) in
\eqref{density:central-derivatives} is now fixed as well.
We give a finite list of requirements on \(Q\), all independent of
the induction stage.

The horizontal derivative estimate at stage \(j-1\), integrated along
the radial arc just described, will imply
\begin{equation}\label{density:central-comparison}
 e^{-K_*D_R/\sqrt Q}
 \le\frac{S_p(\theta)}{W_{j-1}(\xi)}
 \le e^{K_*D_R/\sqrt Q}.
\end{equation}
Indeed the logarithmic derivative along a unit horizontal arc is at most
\(K_*\sqrt{k_{j-1}}\), and the arc length is at most \(D_R/\sqrt{k_j}\).
At the initial stage the ratio is exactly one.
Since \(k_j\ge Q\), put
\[
 H(Q)=\exp\left(K_*D_R/\sqrt Q+2R_0/Q\right).
\]

Choose a threshold \(Q_0\) such that every integer \(Q\ge Q_0\)
satisfies all earlier degree and patch requirements, and also
\begin{equation}\label{density:Q-conditions}
 \begin{gathered}
 Q\ge2,\qquad H(Q)\le2,\qquad
 K_*/Q\le1,\qquad C_1/Q^2\le1,\qquad
 Q^{-1/2}\le c_0/2,\\
 H(Q)(F+BK_*/Q)\le1-2c_0.
 \end{gathered}
\end{equation}
These conditions are compatible: \(H(Q)\to1\) and
\(F=1-4c_0<1-2c_0\).

We now fix any integer \(Q\ge Q_0\) and induct on \(j\).
Assume the claimed estimates through stage \(j-1\); at \(j=1\)
the preceding density is the constant one.
Equation~\eqref{density:central-derivatives}, obtained from earlier
stages only, and the horizontal comparison
\eqref{density:central-comparison} establish the normalized hypotheses
of Lemma~\ref{density:normalized-patch}.  Moreover,
\[
 |A_{j,p}|
 \le e^{2y/k}S_p(F+Br_1)
 \le H(Q)(F+BK_*/Q)W_{j-1}
 \le(1-2c_0)W_{j-1}.
\]
At any point at most one patch correction is nonzero.  Hence
\[
 2c_0W_{j-1}\le W_j\le(2-2c_0)W_{j-1},
\]
which proves positivity and the required multiplicative comparison.

For \(Y=Z/\sqrt{k_j}\), the previous derivative bound gives
\(|YW_{j-1}|\le(K_*/\sqrt Q)W_{j-1}\).
For \(Y=T/k_j\), it gives the smaller coefficient \(K_*/Q\).
Together with Lemma~\ref{density:normalized-patch},
\eqref{density:Q-conditions}, and the choice of \(K_*\), these yield
\[
 |YW_j|\le(C_2+K_*/\sqrt Q)W_{j-1}
 \le K_*c_0W_{j-1}\le K_*W_j.
\]
Outside the patches the correction vanishes.  Its smooth zero
extension at their boundaries gives the same estimate there.
Multiplying by \(\sqrt{k_j}\) or \(k_j\) proves the two derivative
bounds and completes the induction.

Lemma~\ref{density:moments} supplies the exact moments, rotation mean
and bandwidth independently of these inequalities.  In particular
\(\int_SW_j\,d\sigma=1\); positivity therefore makes
\(W_j\,d\sigma\) a probability measure.
\end{proof}

We emphasize the order of the choices.  The normalized correction
constant \(C_2\) depends only on the fixed patch data.  It determines
\(K_*\); this in turn determines \(C_1\); finally \(Q_0\) is chosen.
No lower bound for \(W_j\) uniform in \(j\) is required.  The positive
comparison at each stage suffices for logarithmic differentiation,
and all later integral estimates will use the probability normalization.

\section{Regularized logarithmic tests}\label{sec:tests}

Keep the profile parameters $a,\alpha,s_1,f,b,R_0$ from
Lemma~\ref{peaks:profiles}.  We shall construct a smooth
plurisubharmonic function at each degree whose average against its
representing density is uniformly negative.  The same negative bound
will hold against all later densities.  The estimates are uniform for
all sufficiently large integers $Q$, so the final value of $Q$ can be
chosen after the regularization and spacing parameters.

A positive measure representing holomorphic evaluation need not satisfy
the logarithmic inequality
$\log|f(0)|\le\int_S\log|f|\,d\mu$ for every holomorphic $f$
continuous on the closed ball. Hedenmalm gives an explicit example of
this distinction between representing and Jensen measures
\cite[Theorem~3.7]{Hedenmalm1989}.
Here we need smooth logarithmic tests and negative bounds that persist
through all later density corrections. We prove these properties
directly for the densities of Theorem~\ref{density:main}.

Fix a constant $C_P$ in the bound
$|P_k(z)|\le C_P|z|^k$ of Lemma~\ref{peaks:bounds}.

For $\epsilon>0$, define
\begin{equation}\label{tests:regularized}
 F_\epsilon(v)=
 \log\frac{|1-av|^2+\epsilon^2}{1+\epsilon^2},
 \qquad U_k(z)=F_\epsilon(P_k(z)).
\end{equation}
In particular $F_\epsilon(0)=U_k(0)=0$.  The function $U_k$ is smooth
and plurisubharmonic on $\mathbb C^2$, since
\begin{equation}\label{tests:hessian}
 (U_k)_{i\bar j}
 =\frac{a^2\epsilon^2}
        {(|1-aP_k|^2+\epsilon^2)^2}
   (P_k)_i\overline{(P_k)_j}.
\end{equation}
For fixed $\epsilon$ and $A<\infty$, all real derivatives of
$F_\epsilon$ of any fixed order are bounded on $|v|\le A$.
The mean-value estimate also gives
\begin{equation}\label{tests:lipschitz}
 |F_\epsilon(v)|\le C_{\epsilon,A}|v|
 \qquad (|v|\le A).
\end{equation}
All subsequent regularization constants may depend on the fixed
$\epsilon$.

\begin{lemma}[Exact averaging]\label{tests:averaging}
Let $k=k_j=Q^j$, where $Q$ satisfies the hypotheses of
Theorem~\ref{density:main}.  For $0\le r\le1$,
\begin{equation}\label{tests:averaging-identity}
 \int_S U_k(r\xi)W_j(\xi)\,d\sigma
 =
 \int_S U_k(r\xi)\,d\sigma
 +\sum_p\int_{\mathcal U_{k,p}}
     U_k(r\xi)e^{2y/k}f(y)\cos(k\theta)\,d\sigma.
\end{equation}
\end{lemma}

\begin{proof}
Under the global phase rotation $\xi\mapsto e^{i\omega}\xi$,
homogeneity gives $P_k(e^{i\omega}\xi)=e^{ik\omega}P_k(\xi)$.
Thus $U_k(r\xi)$ has only phase frequencies in $k\mathbb Z$.
By Lemma~\ref{density:moments}, the phase bandwidth of $W_{j-1}$
is $\ell_{j-1}<k$ and its constant coefficient on each orbit is one.
Averaging the product over the phase therefore gives
\[
 \int_S U_k(r\xi)W_{j-1}(\xi)\,d\sigma
 =\int_S U_k(r\xi)\,d\sigma.
\]

Each patch is phase-invariant.  Its radial variable $y$ is fixed
under rotation, while $\theta$ and $\gamma$ both increase by $\omega$.
Write the central restriction as
$S_p(\theta)=W_{j-1}(e^{i\theta}p)
 =\sum_s\widehat S_p(s)e^{is\theta}$.
In the patch increment of~\eqref{density:recursion}, the frequencies
are $\pm(k+s)$ with $|s|\le\ell_{j-1}<k$.  Pairing with a frequency in
$k\mathbb Z$ is possible only for $s=0$.
Since $\widehat S_p(0)=1$ and $S_p'$ has zero constant coefficient,
the surviving part of the increment
is exactly $e^{2y/k}f(y)\cos(k\theta)$.  Integrating this phase average
proves~\eqref{tests:averaging-identity}.
\end{proof}

The next estimate chooses the regularization before the spacing.
This order matters: the single-peak approximation will be applied
to a fixed smooth logarithm, and its error can then be made small by
increasing $D_0$.

\begin{proposition}[Negative averages at one scale]\label{tests:negative}
There exist $\epsilon>0$, $\eta\in(0,1/2)$ and $D_*\ge1$, depending
only on the fixed profile data, with the following property.
For every $D_0\ge D_*$ there are $k_*>0$ and $d_2>0$ such that, for
every integer $Q$ satisfying the density theorem and $Q\ge k_*$,
\begin{equation}\label{tests:one-scale-negative}
 \int_S U_{k_j}(r\xi)W_j(\xi)\,d\sigma\le-d_2
 \qquad
 \bigl(j\ge1,\quad 1-\eta\le r^{k_j}\le1\bigr).
\end{equation}
The constants $k_*,d_2$ may depend on $D_0$, but are independent of
$j$ and of further increases in $Q$.
\end{proposition}

\begin{proof}
Put
\[
 \Delta=2s_1(\cosh\alpha-1)-\alpha^2>0.
\]
We first compute the integral for a single peak:
\begin{equation}\label{tests:model}
 I_\epsilon(\rho)=
 \int_0^{R_0}\int_0^{2\pi}
 \log\frac{|1-a\rho e^{-y}e^{it'}|^2+\epsilon^2}
               {1+\epsilon^2}
 (1+f(y)\cos t')\,\frac{dt'}{2\pi}\,dy.
\end{equation}
At $\rho=1$, the unregularized logarithm has only one singularity,
at $(y,t')=(\alpha,0)$ modulo $2\pi$.  Near that point,
\[
 |1-e^{\alpha-y+it'}|^2
 \asymp (y-\alpha)^2+(t')^2.
\]
Its logarithm is locally integrable.  For $0<\epsilon\le1$, the
regularized logarithm is bounded above uniformly, and its negative
part is bounded by the negative part of the unregularized logarithm
plus $\log2$.  Dominated convergence therefore applies in
\eqref{tests:model}.

For $q<1$, the mean and cosine moment of
$\log|1-qe^{it'}|^2$ are $0$ and $-q$, respectively, by the logarithm
series.  For $q>1$, factoring out $q^2$ gives mean $2\log q$ and
cosine moment $-q^{-1}$.  Using the cancellation
$\int_0^{R_0}e^{-y}f(y)\,dy=0$, we obtain
\begin{align}
 I_0(1)
 &=\alpha^2-\int_0^\alpha f(y)e^{y-\alpha}\,dy
             -\int_\alpha^{R_0}f(y)e^{\alpha-y}\,dy \notag\\
 &=\alpha^2+\int_0^\alpha
       f(y)\bigl(e^{\alpha-y}-e^{-(\alpha-y)}\bigr)\,dy \notag\\
 &=\alpha^2-2s_1(\cosh\alpha-1)=-\Delta.
 \label{tests:model-sign}
\end{align}
The last equality uses $f=-s_1$ on $[0,\alpha]$.

We now compare this model with~\eqref{tests:averaging-identity}.
For $k=k_j$ put $\rho=r^k$, so homogeneity gives
$U_k(r\xi)=F_\epsilon(\rho P_k(\xi))$.
Off the patches, $|aP_k|\le1/2$ by Lemma~\ref{peaks:profiles}.
If $v=\rho P_k(\xi)$, $0\le\rho\le1$, and $s=|1-av|^2$, then
$1/4\le s\le9/4$ and
\[
 F_\epsilon(v)-\log|1-av|^2
 =\log(1+\epsilon^2/s)-\log(1+\epsilon^2).
\]
The derivative in $s$ of the first term has magnitude at most
$16\epsilon^2$, and $|s-1|\le C|v|$.  Hence
\[
 |F_\epsilon(v)-\log|1-av|^2|
 \le C\epsilon^2|P_k(\xi)|.
\]
The off-patch set is phase-invariant, and the phase mean of the
unregularized logarithm is zero.  By the $L^1$ estimate of
Lemma~\ref{peaks:bounds}, its total contribution is thus at most
\begin{equation}\label{tests:off-patch}
 C_{\rm off}\epsilon^2\,\frac{N_k}{k}.
\end{equation}
Here $C_{\rm off}$ is independent of $D_0,k,\rho$ and
$0<\epsilon\le1$.

On a patch, the measure formula of Lemma~\ref{peaks:patches} writes
the combined contribution as
\[
 \frac2k\int_0^{R_0}\int_{\mathbb T^2}
 F_\epsilon(\rho P_k(\xi))
       \bigl(e^{-2y/k}+f(y)\cos(k\theta)\bigr)
       \,\frac{d\theta\,d\gamma}{(2\pi)^2}\,dy.
\]
For fixed $\epsilon,R_0$, replace $P_k(\xi)$ by
$e^{-y}e^{ik\theta}$ using the single-peak estimate, and replace
$e^{-2y/k}$ by one.  The bounded derivative of $F_\epsilon$ and the
bound $|e^{-2y/k}-1|\le2R_0/k$ show that the result differs from
$(2/k)I_\epsilon(\rho)$ by at most
\begin{equation}\label{tests:patch-error}
 \frac2k C_{\epsilon,R_0}
 \left(e^{-c(D_0-\sqrt{2R_0})^2}+\frac1k\right).
\end{equation}
All constants are uniform over patches and $0\le\rho\le1$.
The change from $k\theta$ to $t'$ in the angular average is exact,
as $k$ is an integer.

Choose $\epsilon\in(0,1]$ so small that
\[
 |I_\epsilon(1)+\Delta|\le\Delta/8,
 \qquad C_{\rm off}\epsilon^2\le\Delta/8.
\]
By continuity of~\eqref{tests:model} in $\rho$, choose
$\eta\in(0,1/2)$ such that
$I_\epsilon(\rho)\le-3\Delta/4$ for $1-\eta\le\rho\le1$.
Next choose $D_*$ large enough for patch disjointness and for the
first error in~\eqref{tests:patch-error} to be at most $\Delta/16$
inside its factor $2/k$.  For each $D_0\ge D_*$, take $k_*$ large
enough for all the geometric estimates and for the remaining error
inside this factor to be at most $\Delta/16$.
Summing the patches and adding~\eqref{tests:off-patch} gives, in
particular,
\[
 \int_S U_k(r\xi)W_j\,d\sigma
 \le-\frac{\Delta}{2}\frac{N_k}{k}
 \le-\frac{\Delta}{2D_0^2}.
\]
Thus $d_2=\Delta/(2D_0^2)$ is a valid choice.
\end{proof}

\begin{lemma}[Negative averages against later densities]\label{tests:later}
Fix $\epsilon,\eta,D_0,k_*,d_2$ as in
Proposition~\ref{tests:negative}, and put $d_3=d_2/2$.
There is a threshold $Q_*$ such that, for every integer $Q\ge Q_*$,
the densities satisfy
\begin{equation}\label{tests:later-negative}
 \int_S U_{k_j}(r\xi)W_N(\xi)\,d\sigma\le-d_3
 \qquad
 \bigl(N\ge j\ge1,\quad1-\eta\le r^{k_j}\le1\bigr).
\end{equation}
\end{lemma}

\begin{proof}
Include the density and degree thresholds in $Q_*$, and require
$Q\ge4$.  By Lemma~\ref{density:moments}, the increment
$W_\nu-W_{\nu-1}$ has frequencies within $\ell_{\nu-1}$ of
$\pm k_\nu$.  Since $\ell_{\nu-1}\le2k_{\nu-1}$, each such
frequency has absolute value at least $(Q-2)k_{\nu-1}$.
Summing the increments with $\nu>j$ shows that every frequency in
$W_N-W_j$ has absolute value at least $(Q-2)k_j$.

For fixed $\xi,r$, write the logarithmic profile as
\[
 V(t')=\log\frac{|1-Ae^{it'}|^2+\epsilon^2}{1+\epsilon^2},
 \qquad A=a r^{k_j}P_{k_j}(\xi).
\]
The uniform polynomial bound gives $|A|\le aC_P$.  Differentiating
the logarithm shows that
\[
 \|V''\|_\infty\le
 B_\epsilon:=
 \frac{2aC_P}{\epsilon^2}
 +\frac{4a^2C_P^2}{\epsilon^4}.
\]
Two integrations by parts bound its $m$th Fourier coefficient by
$B_\epsilon/m^2$ for $m\ne0$. To track these coefficients as functions
of $\xi$, define
\[
 c_m(\xi)=\int_0^{2\pi}
 F_\epsilon\bigl(r^{k_j}P_{k_j}(\xi)e^{it'}\bigr)
 e^{-imt'}\,\frac{dt'}{2\pi}.
\]
Homogeneity and a change of variable give
$c_m(e^{i\omega}\xi)=e^{imk_j\omega}c_m(\xi)$.
Thus $c_m$ has global phase frequency $mk_j$, and
$U_{k_j}(r\xi)=\sum_{m\in\mathbb Z}c_m(\xi)$.
The coefficient bound makes this series uniformly convergent.
Its truncation to $|m|\le Q-3$ has uniform error at most
$2B_\epsilon/(Q-3)$ and only frequencies of absolute value less
than $(Q-2)k_j$. Phase averaging therefore makes its integral
against $W_N-W_j$ vanish.
Both densities are positive probabilities by
Theorem~\ref{density:main}; hence
$\|W_N-W_j\|_{L^1(\sigma)}\le2$.  It follows that
\[
 \left|\int_S U_{k_j}(r\xi)(W_N-W_j)\,d\sigma\right|
 \le\frac{4B_\epsilon}{Q-3}.
\]
Choose $Q_*$ so that this is at most $d_2/2$, and apply
\eqref{tests:one-scale-negative}.
\end{proof}

\section{The controlled base potential}\label{sec:base-potential}

We now sum the logarithmic tests to prove
Proposition~\ref{base:existence}.  Fix parameters and an integer $Q$
for which Theorem~\ref{density:main} and Lemma~\ref{tests:later}
hold, and write $k_j=Q^j$.  For a finite amplitude $L>0$, set
\begin{equation}\label{base:potential}
 \phi(z)=\psi(z)+L\sum_{j=1}^\infty U_{k_j}(z).
\end{equation}
We first justify this definition.  We then choose $L$ to exclude
the required pluriharmonic minorant, and finally prove the uniform
chart estimates needed by the metric construction.

\begin{lemma}[Smooth convergence]\label{base:convergence}
For every finite $L>0$, the series in~\eqref{base:potential}
converges locally in $C^\infty(\mathbb B)$.  Its sum is smooth,
satisfies $\phi(0)=0$, and has complex Hessian $h\ge I$.
\end{lemma}

\begin{proof}
Fix $0<r_0<r_1<1$.  Around each point of $|z|\le r_0$, the complex
polydisk of radius $(r_1-r_0)/(2\sqrt2)$ lies in $|z|<r_1$.
Cauchy estimates there, together with Lemma~\ref{peaks:bounds}, give
\[
 \sup_{|z|\le r_0}|D^mP_k(z)|
 \le C_{m,r_0,r_1}r_1^k
\]
for every fixed derivative order $m$.  Here and below $D^m$ denotes
any real partial derivative of total order $m$.
The values of $P_k$ remain in a fixed bounded disk.
Every positive-order derivative of $F_\epsilon(P_k)$ is a sum of
products containing at least one positive-order derivative of
$P_k$ or its conjugate.  The bounded derivatives of $F_\epsilon$
therefore give the same exponential estimate for $D^mU_k$.
For $m=0$, use~\eqref{tests:lipschitz}.
Since $\sum_j r_1^{Q^j}<\infty$, the series converges in every local
$C^m$ norm.

Each summand vanishes at zero and is plurisubharmonic by
\eqref{tests:hessian}.  Termwise differentiation is justified by
the local $C^2$ convergence, and hence
\[
 \partial\bar\partial\phi
 =I+L\sum_j\partial\bar\partial U_{k_j}\ge I.
\]
\end{proof}

\begin{lemma}[Absence of a pluriharmonic minorant]\label{base:minorant}
If
\begin{equation}\label{base:amplitude}
 L>\frac{5\log Q}{d_3},
\end{equation}
there are no $H\in\mathcal O(\mathbb B)$ and $C_H\in\mathbb R$
satisfying~\eqref{base:no-minorant}.
\end{lemma}

\begin{proof}
Put
\begin{equation}\label{base:radii}
 r_N=1-\frac{\eta}{k_N}.
\end{equation}
For every $j\le N$, Bernoulli's inequality gives
\[
 r_N^{k_j}
 =\left(1-\frac{\eta}{k_N}\right)^{k_j}
 \ge1-\eta\frac{k_j}{k_N}\ge1-\eta.
\]
Thus all of the first $N$ modes satisfy the hypothesis of
\eqref{tests:later-negative} at this same radius.
For $j=N+m>N$, on the other hand,
\[
 r_N^{k_j}\le e^{-\eta k_j/k_N}=e^{-\eta Q^m}.
\]
By the uniform polynomial bound and~\eqref{tests:lipschitz},
\begin{equation}\label{base:tail}
 \sum_{j>N}\sup_{\xi\in S}|U_{k_j}(r_N\xi)|
 \le C_\epsilon\sum_{m\ge1}e^{-\eta Q^m}
 =:C_{\rm tail}<\infty,
\end{equation}
independently of $N$.  Because $W_N$ is a positive probability,
the same bound controls the weighted tail, without requiring a
uniform pointwise bound on $W_N$.

The base potential contributes only one logarithm:
\[
 \psi(r_N\xi)=-\log(1-r_N^2)
 \le\log\frac1{1-r_N}=\log(k_N/\eta).
\]
Combining this with~\eqref{tests:later-negative} and
\eqref{base:tail} yields
\begin{align}
 &\int_S\left(\phi(r_N\xi)
                  +4\log\frac1{1-r_N}\right)W_N\,d\sigma \notag\\
 &\hspace{1cm}\le
 5\log(k_N/\eta)+L(-d_3N+C_{\rm tail}) \notag\\
 &\hspace{1cm}=
 N(5\log Q-Ld_3)+5\log(1/\eta)+LC_{\rm tail}
 \longrightarrow-\infty.
 \label{base:negative-means}
\end{align}

By Theorem~\ref{density:main}, the measures
$\mu_N=W_N\,d\sigma$ are positive representing probabilities.
Apply Lemma~\ref{peaks:mean-criterion} to these measures, the radii
$r_N$, and the continuous function
$u(z)=\phi(z)+4\log(1/(1-|z|))$.
The limit~\eqref{base:negative-means} is exactly its hypothesis,
so the required minorant cannot exist.
\end{proof}

It remains to control the metric jets as the chart centers
approach the boundary.  The value of $\phi$ need not be bounded.
For its derivatives, the low degrees become almost constant in a
centered chart, while the high degrees are exponentially small.

\begin{lemma}[Uniform chart estimates]\label{base:chart-estimates}
For every fixed finite $L>0$, the potential~\eqref{base:potential}
has uniformly bounded positive-order derivatives through order four
at the origins of all centered ball charts.  Consequently its
complex Hessian satisfies~\eqref{base:control} and
\eqref{base:jets}, with $c=1$ and a finite constant $C$.
\end{lemma}

\begin{proof}
After a unitary change of coordinates, put the center at $(r,0)$,
where $0\le r<1$, and let $\delta=1-r$.  Rotating $P_k$ preserves its
degree and its uniform homogeneous bound.  Use the chart $T_r$ of
\eqref{base:charts}, on the fixed complex polydisk
\[
 D_\rho=\{|\zeta_1|<\rho,\ |\zeta_2|<\rho\},
 \qquad \rho=\frac1{16}.
\]
The denominators in $T_r$ have modulus at least $1-\rho$.
We estimate the polynomial in the two ranges of $k\delta$.

\paragraph{High degrees: $k\delta>1$.}
The ball automorphism identity gives, on $D_\rho$,
\[
 1-|T_r(\zeta)|^2
 =\frac{(1-r^2)(1-|\zeta|^2)}{|1+r\zeta_1|^2}
 \ge \frac{1-2\rho^2}{(1+\rho)^2}\,\delta.
\]
It follows that $1-|T_r(\zeta)|\ge c\delta$, with a fixed $c>0$.
Thus
\[
 |P_k(T_r(\zeta))|\le C_P e^{-ck\delta}.
\]
Cauchy estimates in this fixed polydisk bound every fixed-order
holomorphic derivative at zero, including order zero, by
$C_m e^{-ck\delta}$.

\paragraph{Low degrees: $k\delta\le1$.}
Consider the polydisk in the original coordinates
\[
 A_k=\left\{|z_1-r|<\frac1k,\quad
                 |z_2|<\frac1{\sqrt k}\right\}.
\]
It may extend outside the ball, but $P_k$ is an entire homogeneous
polynomial.  The bound of Lemma~\ref{peaks:bounds} therefore holds
there as well.  In fact
\[
 |z|^2\le(1+1/k)^2+1/k\le1+4/k,\qquad
 |P_k(z)|\le C_P(1+2/k)^k\le C_Pe^2.
\]
Cauchy estimates on the concentric half-polydisk give
\[
 |\partial_1P_k|\le Ck,\qquad
 |\partial_2P_k|\le C\sqrt k.
\]
Meanwhile the chart displacements satisfy
\[
 |T_r^1(\zeta)-r|
 \le\frac{2\rho}{1-\rho}\delta,\qquad
 |T_r^2(\zeta)|
 \le\frac{\sqrt2\rho}{1-\rho}\sqrt\delta.
\]
Both numerical coefficients are less than $1/2$.  As
$\delta\le1/k$, the chart image lies in the half-polydisk just
used.  Integrating the gradient along the straight segment from
$(r,0)$ to $T_r(\zeta)$ gives
\[
 |P_k(T_r(\zeta))-P_k(r,0)|
 \le C(k\delta+\sqrt{k\delta})
 \le2C\sqrt{k\delta}.
\]
This segment stays in the half-polydisk by convexity.  Apply
Cauchy estimates on $D_\rho$ to the holomorphic difference
$P_k\circ T_r-P_k(r,0)$.  Every positive-order derivative at zero
is bounded by $C_m\sqrt{k\delta}$.  We make no such assertion for
the zeroth-order value in this range.

\paragraph{Composition and summation.}
In both ranges the values of $P_k\circ T_r$ stay in a fixed bounded
disk.  Each positive-order real derivative of
$U_k\circ T_r=F_\epsilon(P_k\circ T_r)$ is a sum of products
containing differentiated factors of $P_k\circ T_r$ or its
conjugate.  For $1\le m\le4$, the preceding bounds and the smooth
composition estimate give
\begin{equation}\label{base:scale-jets}
 |D^m(U_k\circ T_r)(0)|
 \le C_{\epsilon,m}
 \begin{cases}
  \sqrt{k\delta},&k\delta\le1,\\
  e^{-ck\delta},&k\delta>1.
 \end{cases}
\end{equation}
Indeed every term contains at least one differentiated factor,
and a product of several factors obeys the same bound because the
displayed scale factor is at most one.  This proves the estimate
for all real derivatives, hence for all pure and mixed complex
derivatives of these orders.

Let $J$ be the largest index with $Q^J\delta\le1$, or put $J=0$
if there is no such index.  The low-degree contribution is bounded
by the geometric sum
\[
 \sum_{j\le J}\sqrt{Q^j\delta}
 \le\frac1{1-Q^{-1/2}}.
\]
For the remaining degrees, the first value of $Q^j\delta$ exceeds
one, so
\[
 \sum_{j>J}e^{-cQ^j\delta}
 \le\sum_{m\ge0}e^{-cQ^m}<\infty.
\]
Both bounds are uniform in the chart center.  Termwise
differentiation is legitimate by Lemma~\ref{base:convergence}.
Multiplication by the fixed amplitude $L$ preserves these bounds.

Finally, the base potential transforms as
\[
 \psi\circ T_r
 =-\log(1-|\zeta|^2)-\log(1-r^2)
      +\log|1+r\zeta_1|^2.
\]
Its positive-order derivatives through order four at zero are
uniformly bounded.  The last term is pluriharmonic, so the complex
Hessian of $\psi\circ T_r$ is exactly the standard ball Hessian in
these coordinates.  The derivatives of $\phi\circ T_r$ needed for
the pulled-back Hessian components and their first and second
derivatives have orders two, three and four.  The bounds just
proved, together with $h\ge I$ from
Lemma~\ref{base:convergence}, establish the assertion.
\end{proof}

\begin{proof}[Proof of Proposition~\ref{base:existence}]
Choose $a,\alpha,s_1,f,b,R_0$ as in Lemma~\ref{peaks:profiles}.
Proposition~\ref{tests:negative} first chooses $\epsilon,\eta$ and
then permits any sufficiently large spacing $D_0$.  Fix such a
$D_0$ and its degree threshold $k_*$, and fix
$d_3=d_2/2>0$.  The density theorem holds for every $Q$ above a
threshold depending only on the fixed patch data and spacing.
Its derivative constant $K_*$ is fixed before that threshold is
chosen.  We may therefore choose one integer $Q$ satisfying the
density threshold, $Q\ge k_*$, and the Fourier-error threshold of
Lemma~\ref{tests:later}.

Now fix a finite amplitude $L$ satisfying~\eqref{base:amplitude}.
Lemma~\ref{base:convergence} constructs the smooth potential with
$h\ge I$; Lemma~\ref{base:minorant} proves the required obstruction;
and Lemma~\ref{base:chart-estimates} supplies its finite uniform
chart bounds.  The chart constant is allowed to depend on these
fixed choices, including $L$.  The metric parameter $\lambda$ in
Proposition~\ref{metric:transfer} is chosen only afterward from
these bounds, so none of the choices above depends on $\lambda$.
\end{proof}

\section{The global Hartogs metric}\label{sec:hartogs}

We now turn to Proposition~\ref{metric:transfer}.  In this section we
prove all its global geometric assertions.  The real sectional-curvature
bounds will follow from the calculation in Section~\ref{sec:curvature}
and the estimates in Section~\ref{sec:pinching}.

The potential specializes Calabi's construction on a Hermitian line
bundle \cite[Section~3]{Calabi1979}. We verify its positivity,
smoothness, and completeness directly, including at the zero section.

\begin{lemma}\label{metric:complete}
Let $\phi$ be a smooth real plurisubharmonic function on $\B$, and let
$\lambda>0$.  The domain $M_\phi$ in \eqref{metric:domain} is
contractible, and the complex Hessian of
$\mathcal K_\lambda=\lambda\psi-\log(1-e^\phi|w|^2)$ defines a smooth
geodesically complete K\"ahler metric
$g=2\operatorname{Re}\widehat g$ on $M_\phi$.
Moreover, $\mathcal K_\lambda$ is a smooth strictly plurisubharmonic
exhaustion of $M_\phi$.
\end{lemma}

\begin{proof}
\emph{The domain and a smooth positive metric.}
The function $t=e^\phi|w|^2$ is smooth on $\B\times\C$, so
$M_\phi=\{t<1\}$ is open in $\C^3$.  The map
\[
 (z,w)\longmapsto (z,e^{\phi(z)/2}w)
\]
is a smooth real diffeomorphism from $M_\phi$ to $\B\times\mathbb D$,
where $\mathbb D=\{u\in\C:|u|<1\}$.  Its inverse is
$(z,u)\mapsto(z,e^{-\phi(z)/2}u)$.  Thus $M_\phi$ is contractible;
no holomorphic trivialization is needed for this conclusion.

The potential $\mathcal K_\lambda$ is smooth at every point of
$M_\phi$, including its zero section.  Write
\[
 h=\partial\bar\partial\phi,\qquad
 x=\frac{t}{1-t},\qquad G=\lambda I+xh,
 \qquad \vartheta=dw+w\partial\phi.
\]
The base covectors together with $\vartheta$ form a global smooth
complex coframe.  Differentiating the potential gives
\begin{equation}\label{metric:smooth-form}
 \widehat g
 =G+\frac{e^\phi}{(1-t)^2}\,
       \vartheta\otimes\overline{\vartheta}.
\end{equation}
For a complex tangent vector $u=(u_z,u_w)$ this means
\[
 \widehat g(u,u)
 =G(u_z,u_z)
  +\frac{e^\phi}{(1-t)^2}
      |u_w+w\partial\phi(u_z)|^2.
\]
Since $h\ge0$, the first term controls every nonzero horizontal
projection through $G\ge\lambda I$; when $u_z=0$, the second term is
positive unless $u_w=0$.  Thus $\widehat g$ is positive definite.
At the zero section its matrix is particularly simple:
\begin{equation}\label{metric:axis}
 \widehat g\big|_{(z,0)}
 =\lambda I(z)+e^{\phi(z)}\,dw\otimes d\bar w.
\end{equation}
The mixed blocks vanish.  Being a positive complex Hessian of a smooth
real potential, $\widehat g$ defines the stated K\"ahler metric.

Away from the zero section one can write $s=\log t$ and
$q=x(1+x)=t/(1-t)^2$.  Then $\partial s=\vartheta/w$, so
\eqref{metric:smooth-form} also reads
\begin{equation}\label{metric:log-form}
 \widehat g=\lambda I+xh+
             q\,\partial s\otimes\overline{\partial s}.
\end{equation}
We will use this expression for the curvature computation.  Formula
\eqref{metric:smooth-form}, rather than the coordinate $s$, defines the
metric across $w=0$.

\emph{Two functions with bounded gradient.}
For a real smooth function $f$, the convention
$g=2\operatorname{Re}\widehat g$ gives
\begin{equation}\label{metric:gradient-convention}
 |\nabla f|_g^2=2|\partial f|_{\widehat g^{-1}}^2.
\end{equation}
Indeed, for a real vector $X$ with $(1,0)$ part $u$,
$g(X,X)=2\widehat g(u,u)$ and
$df(X)=2\operatorname{Re}\partial f(u)$, and the identity follows
by taking dual norms.

Globally on $M_\phi$,
\[
 \partial[-\log(1-t)]
 =\frac{e^\phi\bar w}{1-t}\,\vartheta.
\]
In the coframe of \eqref{metric:smooth-form}, the Hermitian metric is
block diagonal.  Therefore
\[
 |\partial[-\log(1-t)]|_{\widehat g^{-1}}^2
 =\frac{|e^\phi\bar w/(1-t)|^2}{e^\phi/(1-t)^2}=t.
\]
This calculation includes $w=0$, where both sides are zero.  For the
base potential, direct inversion of its complex Hessian gives
\[
 |\partial\psi|_{I^{-1}}^2=|z|^2.
\]
The covector $\partial\psi$ has only base components in the same
coframe, and $G\ge\lambda I$.  Combining these facts with
\eqref{metric:gradient-convention} yields
\begin{equation}\label{metric:gradient-bounds}
 |\nabla[-\log(1-t)]|_g^2=2t<2,
 \qquad
 |\nabla\psi|_g^2\le\frac{2|z|^2}{\lambda}\le\frac2\lambda.
\end{equation}

\emph{Compact containment of metric balls.}
Fix $p\in M_\phi$ and a finite radius $R\ge0$.  Integrating
\eqref{metric:gradient-bounds} along arbitrary piecewise smooth paths
and taking the infimum of their lengths shows that
$d_g(p,q)\le R$ implies
\[
 -\log(1-t(q))\le a_R:=-\log(1-t(p))+\sqrt2R,
 \qquad
 \psi(q)\le b_R:=\psi(p)+\sqrt{2/\lambda}\,R.
\]
This argument does not require a minimizing geodesic.  Set
\[
 \tau_R=1-e^{-a_R}<1,\qquad
 \rho_R=\sqrt{1-e^{-b_R}}<1.
\]
Every such $q=(z,w)$ satisfies $t(q)\le\tau_R$ and
$|z|\le\rho_R$.  Let $m_R$ be the minimum of $\phi$ on the compact
base ball $\{|z|\le\rho_R\}$.  Then
\[
 |w|^2=e^{-\phi(z)}t(q)\le e^{-m_R}\tau_R.
\]
Consequently the metric ball is contained in
\begin{equation}\label{metric:compact-containment}
 K_R=\{(z,w):|z|\le\rho_R,
                  \ e^{\phi(z)}|w|^2\le\tau_R\}.
\end{equation}
This set is closed and bounded in $\C^3$, hence compact, and the
strict bounds $\rho_R<1$ and $\tau_R<1$ place it entirely inside
$M_\phi$.  In particular, a bounded-length curve cannot escape through
the base boundary, through a fiber boundary, or through unbounded fiber
coordinates.  The zero section is an interior locus with the smooth
positive metric \eqref{metric:axis}.

The same compact-containment argument applies to sublevel sets of
$\mathcal K_\lambda$.  Both $\psi$ and $-\log(1-t)$ are nonnegative,
so $\mathcal K_\lambda\le R$ implies
$\psi\le R/\lambda$ and $-\log(1-t)\le R$.  These bounds place the
sublevel set in a compact set of the form
\eqref{metric:compact-containment}; the sublevel set is closed there,
hence compact.  Positivity of its complex Hessian was proved above,
so $\mathcal K_\lambda$ is a strictly plurisubharmonic exhaustion.

\emph{Geodesic completeness.}
Suppose a unit-speed geodesic has a finite maximal forward time $T$.
Its image lies in the compact set $K_T$.  On this set the smooth
positive metric has a uniform positive Euclidean lower bound, so the
geodesic's Euclidean velocity is bounded.  The Christoffel symbols are
also bounded on $K_T$, and the geodesic equation then bounds its
Euclidean acceleration.  Both position and velocity have limits as
time tends to $T$, with limiting position in $K_T\subset M_\phi$.
Local existence for the smooth geodesic equation extends the geodesic
beyond $T$, a contradiction.  The same argument in reverse time proves
geodesic completeness.
\end{proof}

The hypotheses of Proposition~\ref{metric:transfer} imply $h\ge cI>0$,
so Lemma~\ref{metric:complete} applies for every $\lambda>0$ under
consideration.  It remains to choose $\lambda$ large enough for the two
uniform real-sectional bounds.  We do this in the next two sections,
keeping the convention $g=2\operatorname{Re}\widehat g$ throughout.

\section{Curvature on real two-planes}\label{sec:curvature}

We retain the base potential and Hartogs metric of
Proposition~\ref{metric:transfer}.  Section~\ref{sec:pinching} will prove
the two uniform sectional-curvature bounds.  Here we compute the full
curvature tensor and reduce its value on an arbitrary real two-plane
to a scalar expression.  In particular, the calculation includes real
planes that are not complex lines.

Throughout this section, constants in tensor estimates depend only on
the uniform centered-chart bounds for $h=(\phi_{i\bar j})$ in
Proposition~\ref{metric:transfer}; the base dimension is two.
We take $\lambda\ge1$.  The base form remains
$I=(\psi_{i\bar j})$, with $\psi=-\log(1-|z|^2)$.

\subsection{A pointwise holomorphic gauge and the tensor blocks}

Fix a point with $w\ne0$, choose centered ball coordinates at its base
point, and choose a local branch of $\log w$.  Put
\[
 s=\log t=\phi+\log w+\overline{\log w},\qquad
 x=\frac{t}{1-t},\qquad q=x(1+x).
\]
In these coordinates the centered base point is zero.  Remove the first
and pure second derivatives of the weight by the holomorphic fiber change
\[
 P(z)=\sum_i\phi_i(0)z_i+
       \frac12\sum_{i,k}\phi_{ik}(0)z_i z_k,
 \qquad v=\log w+P(z).
\]
More precisely, $s=v+\bar v+\widetilde\phi$, where
$\widetilde\phi=\phi-P-\bar P$.  At the point under consideration,
\begin{equation}\label{curv:gauge}
 s_i=s_{ik}=0,\qquad s_v=s_{\bar v}=1,\qquad
 s_{i\bar j}=h_{i\bar j}.
\end{equation}
All higher derivatives of $s$ involving a vertical index vanish.
Subtracting $P+\bar P$ changes neither $h$ nor its component
derivatives in the centered base coordinates.

The metric at the point consequently has blocks
\[
 \widehat g=\begin{pmatrix}G&0\\0&q\end{pmatrix},
 \qquad G=\lambda I+xh.
\]
At the center, $I$ is the identity matrix.  Normalize the vertical
vector to $q^{-1/2}\partial_v$ and denote its index by $0$.
Horizontal indices $i,j,k,l$ range over $1,2$.
In this pointwise frame the metric is $G+V$, where $V$ is the
Hermitian form whose only nonzero entry is $V_{0\bar0}=1$.
Horizontal forms are extended by zero on the vertical direction.
For any Hermitian form $P$, define
\begin{equation}\label{curv:symmetric-tensor}
 (S_P)_{a\bar b c\bar d}
   =P_{a\bar b}P_{c\bar d}+P_{a\bar d}P_{c\bar b}.
\end{equation}

Our complex curvature components use the convention
\begin{equation}\label{curv:coordinate-formula}
 R_{a\bar b c\bar d}
 =-\mathcal K_{ac\bar b\bar d}
   +\widehat g^{p\bar e}
       \mathcal K_{ac\bar e}\mathcal K_{p\bar b\bar d}.
\end{equation}
Repeated indices in this formula are summed over all coordinates;
repeated horizontal indices below are summed over $1,2$.
The inverse coefficients satisfy
$\sum_p\widehat g^{p\bar e}\widehat g_{p\bar j}
=\delta_{ej}$, so their numerical array is the transpose of the
ordinary inverse of the displayed metric matrix.
We relate \eqref{curv:coordinate-formula} to real sectional curvature
below, with the fixed convention $g=2\operatorname{Re}\widehat g$.

\begin{proposition}\label{curv:tensor}
In the frame just described, define the horizontal Hermitian form
\[
 D=(1+2x)h-qhG^{-1}h.
\]
The full curvature tensor is
\begin{equation}\label{curv:tensor-identity}
 R=-\lambda S_I-qS_h-S_V
    -(S_{D+V}-S_D-S_V)+T_4+T_3,
\end{equation}
where $T_j$ has exactly $j$ horizontal indices.  Its nonzero
components, up to the K\"ahler symmetries and conjugation, are
\begin{align}
 (T_4)_{i\bar j k\bar l}
   &=-x h_{i\bar j,k\bar l}
      +x^2G^{p\bar e}h_{i\bar e,k}h_{p\bar j,\bar l},
       \label{curv:four-horizontal-error}\\
 (T_3)_{i\bar j k\bar0}
   &=\sqrt q\bigl(-h_{i\bar j,k}
           +xG^{p\bar e}h_{i\bar e,k}h_{p\bar j}\bigr).
       \label{curv:three-horizontal-error}
\end{align}
Measured by the Euclidean component norms in this frame, uniformly
for $\lambda\ge1$,
\begin{equation}\label{curv:component-errors}
 \|T_4\|\le C_4x,\qquad \|T_3\|\le C_3\sqrt q.
\end{equation}
If $\lambda$ is at least the uniform upper bound for the eigenvalues
of $h$ at centered-chart origins, then
\begin{equation}\label{curv:D-bounds}
 (1+x)h\le D\le(1+2x)h.
\end{equation}
\end{proposition}

\begin{proof}
Write $F(s)=-\log(1-e^s)$.  Direct differentiation gives
\[
 F'=x,\qquad F''=q,\qquad F'''=q(1+2x),\qquad
 F''''=q(1+6x+6x^2).
\]
The centered-chart expression for $\psi$ differs from
$-\log(1-|z|^2)$ by a pluriharmonic function and a constant.
Its derivatives of type $(2,1)$ therefore vanish at the center,
and its derivatives of type $(2,2)$ are $S_I$.

For brevity in the derivative tables, put
\[
 A_{ik\bar j}=h_{i\bar j,k},\qquad
 B_{ik\bar j\bar l}=h_{i\bar j,k\bar l}.
\]
Before normalizing the vertical vector, the only possibly nonzero
third derivatives of type $(2,1)$, up to symmetry, are
\begin{equation}\label{curv:third-derivatives}
 \begin{aligned}
 \mathcal K_{ik\bar j}&=xA_{ik\bar j},\qquad
 \mathcal K_{iv\bar j}=qh_{i\bar j},\qquad
 \mathcal K_{vv\bar v}=q(1+2x).
 \end{aligned}
\end{equation}
The fourth derivatives of type $(2,2)$, up to symmetry and
conjugation, are
\begin{equation}\label{curv:fourth-derivatives}
 \begin{aligned}
 \mathcal K_{ik\bar j\bar l}
   &=\lambda(S_I)_{i\bar j k\bar l}
       +xB_{ik\bar j\bar l}+q(S_h)_{i\bar j k\bar l},\\
 \mathcal K_{ik\bar j\bar v}&=qA_{ik\bar j},\\
 \mathcal K_{iv\bar j\bar v}&=q(1+2x)h_{i\bar j},\\
 \mathcal K_{vv\bar v\bar v}&=q(1+6x+6x^2).
 \end{aligned}
\end{equation}
To check that these lists are complete, apply the chain rule to
$F(s)$ and use \eqref{curv:gauge}.  In the four-horizontal
derivative, the two pairings of mixed second derivatives of $s$
give $qS_h$; the third pairing contains
$s_{ik}s_{\bar j\bar l}=0$.  All other terms contain a vanished
first or pure second horizontal derivative.  In particular,
$\mathcal K_{ik\bar v}=\mathcal K_{vv\bar j}=0$,
$\mathcal K_{ik\bar v\bar v}=0$, and the fourth derivatives
with one horizontal and three vertical indices vanish.
No vanishing of pure third horizontal derivatives is needed.

Substitution into \eqref{curv:coordinate-formula} gives
\eqref{curv:four-horizontal-error} and
\eqref{curv:three-horizontal-error}.  For clarity, all remaining
curvature types in the normalized frame are
\begin{equation}\label{curv:remaining-blocks}
 \begin{aligned}
 R_{i\bar j0\bar0}
   &=-(1+2x)h_{i\bar j}
       +q(hG^{-1}h)_{i\bar j}=-D_{i\bar j},\\
 R_{i\bar0k\bar0}&=0,\qquad
 R_{i\bar00\bar0}=0,\qquad
 R_{0\bar00\bar0}=-2.
 \end{aligned}
\end{equation}
The last identity follows because its unnormalized value is
\[
 -q(1+6x+6x^2)+q(1+2x)^2=-2q^2,
\]
and the four vertical normalization factors multiply to $q^{-2}$.
The tensors $S_V$ and $S_{D+V}-S_D-S_V$ give precisely the
vertical and balanced two-horizontal blocks in
\eqref{curv:remaining-blocks}.  The K\"ahler symmetries and
conjugation now account for every component, proving
\eqref{curv:tensor-identity}.

The lower bound $h\ge cI$ implies
\[
 \|G^{-1}\|\le\frac1{\lambda+cx},\qquad
 x^2\|G^{-1}\|\le\frac{x}{c}.
\]
The assumed component-derivative bounds for $h$, applied to
\eqref{curv:four-horizontal-error} and
\eqref{curv:three-horizontal-error}, prove
\eqref{curv:component-errors} with constants independent of
$\lambda\ge1$ and $x>0$.

Finally, $G=\lambda I+xh$ commutes with $h$ at the center.
On an eigendirection of $h$ with eigenvalue $d$, the ratio of $D$
to $h$ is
\[
 1+2x-\frac{x(1+x)d}{\lambda+xd}.
\]
The subtracted fraction is nonnegative, and it is at most $x$
when $\lambda\ge d$.  This proves \eqref{curv:D-bounds}.
\end{proof}

Although the logarithmic fiber coordinate was used only for $w\ne0$,
the tensor at the zero section can be checked directly in regular
coordinates.  A holomorphic fiber change $W=e^{P(z)+a}w$, with
constant $a$ and a polynomial of degree at most two, makes the new
weight have value zero and vanishing first and pure second
derivatives at the base point.  Its mixed Hessian remains $h$.
Through total degree four, the terms relevant to the curvature are
\[
 \lambda\psi+
 |W|^2\left(1+\sum_{i,j}h_{i\bar j}z_i\bar z_j\right)
 +\frac12|W|^4.
\]
At the origin the metric is $\operatorname{diag}(\lambda I,1)$,
all relevant third derivatives vanish, and direct differentiation
gives
\begin{equation}\label{curv:zero-section}
 R_{i\bar j k\bar l}=-\lambda(S_I)_{i\bar j k\bar l},
 \qquad R_{i\bar j0\bar0}=-h_{i\bar j},
 \qquad R_{0\bar00\bar0}=-2,
\end{equation}
with all other types zero.  In the original fiber coordinate the
normalized vertical vector is $e^{-\phi/2}\partial_w$.
Thus the formulas above hold at the zero section as well by setting
$x=q=0$, $G=\lambda I$, $D=h$, and $T_3=T_4=0$, and using this
regular normalized vertical vector.

\subsection{Real curvature and Gram determinants}

Let $U_1,U_2$ be real linearly independent tangent vectors and let
$u_1,u_2$ be their $(1,0)$ parts, so $U_j=u_j+\bar u_j$.
For a Hermitian form $P$, linear in its first argument, define
\begin{equation}\label{curv:gram-quantities}
 \begin{aligned}
 a_P&=P(u_1,u_1)P(u_2,u_2)
             -(\operatorname{Re}P(u_1,u_2))^2,\\
 b_P&=\operatorname{Im}P(u_1,u_2),\qquad
 m_P=a_{P+V}-a_P-a_V\quad\text{when $P$ is horizontal}.
 \end{aligned}
\end{equation}
Thus $a_P$ is the determinant of the real Gram matrix for
$\operatorname{Re}P$, $b_P$ is the imaginary Hermitian pairing, and
$m_P$ is the cross term in that determinant after adding $V$.
For the metric form $\widehat g$, the real squared area is
$4a_{\widehat g}$ because $g=2\operatorname{Re}\widehat g$.
For a tensor $T$ with the K\"ahler symmetries, also put
\[
 \mathcal H_T=T_{u_1\bar u_1u_2\bar u_2}
       -\operatorname{Re}T_{u_1\bar u_2u_1\bar u_2}.
\]
Here contraction with complex vectors has its usual multilinear
meaning.  These definitions allow the horizontal or vertical projection
of the real plane to be degenerate.

\begin{lemma}\label{curv:real-plane}
For the metric convention $g=2\operatorname{Re}\widehat g$ and
the curvature convention \eqref{curv:coordinate-formula},
\begin{equation}\label{curv:sectional}
 K_g(\operatorname{span}_{\mathbb R}\{U_1,U_2\})
       =\frac{\mathcal H_R}{2a_{\widehat g}}.
\end{equation}
Define the matrix
\begin{equation}\label{curv:xi-definition}
 \Xi^{a\bar b}=u_1^a\overline{u_2^b}
                   -u_2^a\overline{u_1^b}.
\end{equation}
For every tensor $T$ with the K\"ahler symmetries and Hermitian
conjugation symmetry,
\begin{equation}\label{curv:xi-contraction}
 \mathcal H_T=-\frac12
 T_{a\bar b c\bar d}\Xi^{a\bar b}\Xi^{c\bar d}.
\end{equation}
For every Hermitian form $P$,
\begin{equation}\label{curv:S-contraction}
 \mathcal H_{S_P}=a_P+3b_P^2.
\end{equation}
\end{lemma}

\begin{proof}
The real squared area of the pair is $4a_{\widehat g}$, because
$g(U_i,U_j)=2\operatorname{Re}\widehat g(u_i,u_j)$.
In particular $a_{\widehat g}>0$ for a real linearly independent
pair.  Write
\[
 A=R_{u_1\bar u_1u_2\bar u_2},\qquad
 B=R_{u_1\bar u_2u_1\bar u_2}.
\]
The K\"ahler symmetries in the expansion on
$U_j=u_j+\bar u_j$ give the real numerator
\[
 R^{\mathbb R}(U_1,U_2,U_2,U_1)
      =A-B-\bar B+A=2\mathcal H_R.
\]
Division by the real squared area proves \eqref{curv:sectional}.
Expanding the right-hand side of \eqref{curv:xi-contraction}
gives the same expression $A-\operatorname{Re}B$ for $T$.
Finally, if $c_P=P(u_1,u_2)$, direct contraction of
\eqref{curv:symmetric-tensor} gives
\[
 \mathcal H_{S_P}
   =P(u_1,u_1)P(u_2,u_2)+|c_P|^2-2\operatorname{Re}(c_P^2)
   =a_P+3b_P^2.
\]
\end{proof}

For example, the disk potential $-\log(1-|w|^2)$ has normalized
complex curvature component $-2$ by
\eqref{curv:remaining-blocks}.  Taking $u_2=iu_1$ in
\eqref{curv:sectional} gives real sectional curvature $-2$.
This fixes the normalization also on complex lines.

\begin{lemma}\label{curv:gram}
For a positive semidefinite Hermitian form $P$,
$a_P\ge b_P^2$.  For the vertical form $V$ one has $a_V=b_V^2$.
If $P$ is horizontal and positive semidefinite, $m_P$ is nonnegative,
linear and monotone in $P$, and
\begin{equation}\label{curv:mixed-gram}
 m_P\ge2\sqrt{a_Pa_V}.
\end{equation}
When $a_Pa_V>0$, equality in \eqref{curv:mixed-gram} holds exactly
when the two real Gram matrices of the pair for $P$ and $V$ are
proportional.
\end{lemma}

\begin{proof}
Complex Cauchy--Schwarz gives
$P(u_1,u_1)P(u_2,u_2)\ge|P(u_1,u_2)|^2$, proving
$a_P\ge b_P^2$.  Equality holds for $V$ because its complex rank
is one.
Let $A$ and $B$ be the two-by-two real Gram matrices with entries
$\operatorname{Re}P(u_i,u_j)$ and
$\operatorname{Re}V(u_i,u_j)$, respectively.  Then
\[
 a_P=\det A,\qquad a_V=\det B,\qquad
 m_P=\operatorname{tr}(\operatorname{adj}(A)B).
\]
For positive definite $A$, put $C=A^{-1/2}BA^{-1/2}$.
The last expression equals $(\det A)\operatorname{tr}C$,
whereas $\sqrt{a_Pa_V}=(\det A)\sqrt{\det C}$.
The inequality $\operatorname{tr}C\ge2\sqrt{\det C}$ proves
\eqref{curv:mixed-gram}, and continuity gives the semidefinite
case.  If both determinants are positive, equality means that
$C$ is a scalar matrix, which is exactly the claimed proportionality.
The displayed formula for $m_P$ is linear in the entries of $P$.
Its nonnegativity for positive semidefinite $P$ then also proves
monotonicity.
\end{proof}

\subsection{Error terms and the curvature numerator}

Decompose $u_j=(u_{j,H},u_{j,0})$ in the pointwise frame.
Write $\Xi_{HH}$ for the horizontal two-by-two block of $\Xi$
and $\Xi_{H\bar0}$ for its horizontal-to-vertical column.
The other mixed block is the negative conjugate transpose of this
column.  Write $\Xi_{\mathrm{mix}}$ for the matrix consisting of
both mixed blocks, with all other entries zero.
All norms of these blocks below are Euclidean Frobenius norms.

Direct expansion of \eqref{curv:xi-definition} gives
\begin{equation}\label{curv:xi-norms}
 \begin{aligned}
 \|\Xi_{HH}\|^2&=2(a_I+b_I^2)\le4a_I,\\
 \|\Xi_{H\bar0}\|^2&=m_I+2b_Ib_V,\qquad
 \|\Xi_{\mathrm{mix}}\|^2=2(m_I+2b_Ib_V)\le4m_I.
 \end{aligned}
\end{equation}
For the second identity, explicitly expand
$\|u_{1,H}\overline{u_{2,0}}
      -u_{2,H}\overline{u_{1,0}}\|^2$.
The last inequality follows from
$|b_Ib_V|\le\sqrt{a_Ia_V}\le m_I/2$,
by Lemma~\ref{curv:gram}.

\begin{lemma}\label{curv:errors}
For every real linearly independent tangent pair,
\begin{equation}\label{curv:contracted-errors}
 |\mathcal H_{T_4}|\le Cx a_I,\qquad
 |\mathcal H_{T_3}|\le C\sqrt q\,\sqrt{a_I m_I}.
\end{equation}
More explicitly, the three-horizontal error satisfies
\begin{equation}\label{curv:mixed-error}
 |\mathcal H_{T_3}|
   \le C\sqrt q\,\|\Xi_{HH}\|\,
                         \|\Xi_{\mathrm{mix}}\|.
\end{equation}
\end{lemma}

\begin{proof}
In \eqref{curv:xi-contraction}, the tensor $T_4$ pairs two
horizontal blocks of $\Xi$.  The tensor $T_3$ pairs one horizontal
block and one mixed block, since it has exactly three horizontal
indices.  Cauchy--Schwarz and \eqref{curv:component-errors} give
\eqref{curv:mixed-error} and
$|\mathcal H_{T_4}|\le Cx\|\Xi_{HH}\|^2$.
Now use \eqref{curv:xi-norms}.  These estimates remain valid when
one of the projected real areas vanishes.
\end{proof}

We can now express the entire real curvature numerator.
By \eqref{curv:S-contraction} and $a_V=b_V^2$,
\[
 \mathcal H_{S_V}=4a_V,\qquad
 \mathcal H_{S_{D+V}-S_D-S_V}=m_D+6b_Db_V.
\]
Consequently \eqref{curv:tensor-identity} gives
\begin{equation}\label{curv:negative-numerator}
 \begin{aligned}
 -\mathcal H_R={}&\lambda(a_I+3b_I^2)+q(a_h+3b_h^2)
        +4a_V+m_D+6b_Db_V\\
       &-\mathcal H_{T_4}-\mathcal H_{T_3}.
 \end{aligned}
\end{equation}
Its denominator in \eqref{curv:sectional} is positive for every real
two-plane.  It remains to prove that
\eqref{curv:negative-numerator} is uniformly positive relative to
that denominator for sufficiently large fixed $\lambda$, and that
the quotient is also uniformly bounded above.  The mixed term
$6b_Db_V$ can have either sign; Section~\ref{sec:pinching} treats
it together with the possible degeneration of the projected planes.

\section{Uniform pinching on real two-planes}\label{sec:pinching}

We now prove the curvature bounds in Proposition~\ref{metric:transfer}.
The identity of Section~\ref{sec:curvature} reduces the problem to
finite-dimensional algebra.  The limiting main expression is nonnegative,
but it can vanish on certain real planes.
We identify those planes and show that the error with three horizontal
indices vanishes to the additional order needed there.

Throughout this section, horizontal forms are extended by zero on the
vertical factor of \(\mathbb C^2\oplus\mathbb C\).
We retain \(I,V,S_P,a_P,b_P,m_P,\Xi\) from
Section~\ref{sec:curvature}.  Unless otherwise indicated, constants depend
only on the fixed bounds \(cI\le h\le CI\) and the component bounds for
the error tensors.

\subsection{The mixed real-plane algebra}

The leading expression below can vanish even when both projections
of the real plane are nonzero.  For an explicit algebraic model, take
$h=I$, $E=2I$, and
$u_1=(e_1,1)$, $u_2=(ie_1,-i)$, where $e_1=(1,0)\in\mathbb C^2$.
Direct substitution gives
\[
 a_h+3b_h^2+4a_V+m_E+6b_Eb_V=4+4+4-12=0,
 \qquad \Xi_{\mathrm{mix}}=0.
\]
The horizontal and vertical projections have opposite complex
orientations.  The next lemma proves nonnegativity for all planes in
the stated range of $E$ and shows that every zero has this mixed-block
vanishing.

\begin{lemma}\label{pinch:algebra}
Let \(h\) be a positive definite horizontal Hermitian form, let \(E\) be
another horizontal Hermitian form with \(E\ge h\), and let \(u_1,u_2\)
be real-linearly independent vectors in \(\mathbb C^2\oplus\mathbb C\).
Define
\[
 P_h(E)=a_h+3b_h^2+4a_V+m_E+6b_Eb_V.
\]
If \(a_ha_V=0\), then \(P_h(E)>0\).
If \(a_ha_V>0\), put \(\tau=|b_h|/\sqrt{a_h}\).
For every real \(\beta\) such that \(E\le(2-\beta)h\),
\begin{equation}\label{pinch:square}
 P_h(E)\ge
 \left(\sqrt{a_h+3b_h^2}-2\sqrt{a_V}\right)^2
 +2\beta(1+\tau)\sqrt{a_ha_V}.
\end{equation}
In particular, \(h\le E\le2h\) implies \(P_h(E)\ge0\).
If equality holds in this last conclusion, then
\[
 a_h=a_V>0,\qquad |b_h|=\sqrt{a_h},\qquad b_Eb_V<0,
\]
and both mixed blocks \(\Xi^{i\bar0}\) and \(\Xi^{0\bar i}\) vanish.
\end{lemma}

\begin{proof}
If \(a_V=0\), then \(b_V=0\), because the vertical factor has complex
dimension one.
If \(a_h=0\), the two horizontal projections are linearly dependent
over \(\mathbb R\), since \(h\) is positive definite.
Their pairing under every Hermitian form is therefore real, so
\(b_E=0\).
Thus \(a_ha_V=0\) implies \(b_Eb_V=0\), and
\[
 P_h(E)\ge a_h+4a_V+m_h\ge a_{h+V}>0.
\]
Here the last inequality uses the positive definiteness of \(h+V\)
and the real linear independence of the pair.

Suppose now that \(A=a_h>0\) and \(U=a_V>0\).
The Hermitian Gram determinant gives \(0\le\tau\le1\).
Make a real change of basis in the plane so that its horizontal
projections, denoted \(u_H,v_H\), are real-orthonormal for \(h\).
Both ratios \(|b_h|/\sqrt{a_h}\) and \(|b_E|/\sqrt{a_h}\) are unchanged
by this operation.
The squared \(h\)-norms of \(u_H+i v_H\) and \(u_H-i v_H\) are
\(2(1+\tau)\) and \(2(1-\tau)\), in some order.
The absolute difference of their squared \(E\)-norms is \(4|b_E|\)
in the new basis.  Comparing one norm from above and the other from
below gives
\begin{equation}\label{pinch:imaginary-bound}
 \frac{|b_E|}{\sqrt A}
 \le\frac{(2-\beta)(1+\tau)-(1-\tau)}2
 =\frac{1+3\tau}{2}-\frac{\beta(1+\tau)}2
 \le\sqrt{1+3\tau^2}-\frac{\beta(1+\tau)}2.
\end{equation}
For either order of the two norms, the displayed right side is an
upper bound, since \(2-\beta\ge1\).
The last inequality follows from
\[
 1+3\tau^2-\left(\frac{1+3\tau}{2}\right)^2
 =\frac34(1-\tau)^2;
\]
it is strict unless \(\tau=1\).
This argument does not require \(\beta\ge0\).

The mixed Gram inequality \eqref{curv:mixed-gram} and the Hermitian
Gram inequality yield
\[
 m_E\ge2\sqrt{a_EU}\ge2|b_E|\sqrt U.
\]
Since \(U=b_V^2\), it follows that
\[
\begin{aligned}
 P_h(E)
 &\ge A+3b_h^2+4U-4|b_E|\sqrt U\\
 &\ge
 \left(\sqrt{A+3b_h^2}-2\sqrt U\right)^2
 +2\beta(1+\tau)\sqrt{AU},
\end{aligned}
\]
which proves \eqref{pinch:square}.

It remains to examine equality when \(h\le E\le2h\).
The zero-area case already treated shows that \(A,U>0\).
Set \(s=\sqrt{A+3b_h^2}\) and \(B=|b_E|\).
The inequalities above, with \(\beta=0\), give
\[
 P_h(E)\ge(s-2\sqrt U)^2+4\sqrt U(s-B)\ge0.
\]
Equality forces \(s=2\sqrt U\), \(B=s\), and equality in the bounds
used to reach this expression.
The strictness in \eqref{pinch:imaginary-bound} gives \(\tau=1\).
Consequently
\[
 U=A,\qquad B=2\sqrt A,\qquad b_Eb_V<0,\qquad
 m_E=2B\sqrt U.
\]
Since \(a_E\ge B^2\) and \(m_E\ge2\sqrt{a_EU}\), these identities force
\(a_E=B^2\) and equality in the mixed Gram inequality.
The horizontal real \(E\)-Gram matrix and the vertical real Gram matrix
are therefore proportional.

In the real \(h\)-orthonormalized basis used above, \(\tau=1\) says
that the horizontal projections have the form
\[
 u_H,\quad \sigma i u_H,\qquad \sigma\in\{1,-1\}.
\]
Their real \(E\)-Gram matrix is a positive scalar multiple of the
identity.  Proportionality implies that the vertical projections have
equal positive norms and real inner product zero.
The sign \(b_Eb_V<0\) selects the opposite complex orientation, so these
projections have the form
\[
 y,\quad-\sigma i y.
\]
Direct substitution now gives
\[
 \Xi^{i\bar0}
 =u_H^i\overline{-\sigma i y}-(\sigma i u_H^i)\bar y=0.
\]
The other mixed block vanishes as well.
Under any real change of basis in the plane, \(\Xi\) is multiplied by
the determinant of that change.  The vanishing therefore holds in the
original basis.
\end{proof}

\subsection{Uniform negativity for large \texorpdfstring{\(\lambda\)}{lambda}}

The next statement allows every error tensor with the specified bounds
and component types.  This uniformity will also give the compactness
argument needed after \(\lambda\) has been fixed.

\begin{lemma}\label{pinch:positivity}
Fix \(0<c\le C<\infty\) and \(0\le C_{\mathrm e}<\infty\).
There is a finite \(\lambda_0=\lambda_0(c,C,C_{\mathrm e})>0\) with the following
property.
Let \(\lambda\ge\lambda_0\), \(x\ge0\), and let \(h\) be a horizontal
Hermitian form with \(cI\le h\le CI\).  Put
\[
 q=x(1+x),\qquad G=\lambda I+xh,\qquad
 D=(1+2x)h-qhG^{-1}h.
\]
Let \(T_4,T_3\) be tensors with the Kähler symmetries and reality
condition, supported respectively on components with four and with
three horizontal indices, and satisfying
\[
 \|T_4\|\le C_{\mathrm e}x,\qquad
 \|T_3\|\le C_{\mathrm e}\sqrt q.
\]
Define
\[
 R=-\lambda S_I-qS_h-S_V-(S_{D+V}-S_D-S_V)+T_4+T_3.
\]
Then \(\mathcal H_R<0\) on every real two-plane.
\end{lemma}

\begin{proof}
Increase the prospective threshold so that \(\lambda\ge\max(1,C)\).
Then \eqref{curv:D-bounds} gives
\[
 (1+x)h\le D\le(1+2x)h.
\]
If the lemma were false, there would be a sequence
\(\lambda\to\infty\) of permitted data and real two-planes for which
\(F=-\mathcal H_R\le0\).
We suppress the sequence index.
By \eqref{curv:negative-numerator},
\begin{equation}\label{pinch:leading-numerator}
 F=\lambda(a_I+3b_I^2)+q(a_h+3b_h^2)+4a_V
 +m_D+6b_Db_V-\mathcal H_{T_4}-\mathcal H_{T_3}.
\end{equation}
The contracted estimates \eqref{curv:contracted-errors} apply to the
allowed error tensors by their component types and symmetries.

\paragraph{The regime \(q/\lambda\to0\).}
Suppose first that a subsequence has this property.
Monotonicity and linearity of the mixed Gram term give
\(m_D\ge c(1+x)m_I\), while
\[
 |b_D|\le\sqrt{a_D}\le C'(1+x)\sqrt{a_I}.
\]
Discarding nonnegative terms in \eqref{pinch:leading-numerator}, we get
\[
 F\ge\lambda a_I+4a_V+c(1+x)m_I
 -C'\bigl((1+x)\sqrt{a_Ia_V}+x a_I+
             \sqrt q\,\sqrt{a_Im_I}\bigr).
\]
Absorb the first product into \(2a_V\), and the last into
\(\frac c2(1+x)m_I\).  Since \(q/(1+x)=x\), this gives
\[
 F\ge[\lambda-C''((1+x)^2+x)]a_I+
 2a_V+\frac c2(1+x)m_I.
\]
Here \((1+x)^2+x=O(1+q)=o(\lambda)\), because \(\lambda\to\infty\).
Every coefficient is eventually positive, whereas
\(a_I+a_V+m_I=a_{I+V}>0\).  This contradicts \(F\le0\).

\paragraph{The regime with \(\lambda/q\) bounded.}
If the preceding regime has no subsequence, we may pass to one on
which \(q/\lambda\) is bounded below by a positive constant.
Then \(\lambda/q\) is bounded and \(x\to\infty\).
The term \(qS_h\) has four horizontal slots, whereas the vertical
term \(S_V\) already has coefficient one.  To balance these two blocks,
multiply horizontal vector components by \(q^{1/4}\).
In the resulting plane choose a real-orthonormal pair for
\(\operatorname{Re}(I+V)\).
The inverse transformation and a real basis change give a basis of
the original plane; the latter change multiplies \(F\) by a positive
factor.
Thus its sign is unchanged.
The rescaled pairs belong to a fixed compact Stiefel manifold, and
their quantities satisfy \(a_{I+V}=1\).

In the rescaled variables, the leading part of
\eqref{pinch:leading-numerator} is
\begin{equation}\label{pinch:rescaled-leading}
 L_x=\frac{\lambda}{q}(a_I+3b_I^2)+P_h(E),
 \qquad E=\frac D{\sqrt q}.
\end{equation}
The transformed tensors, denoted \(\widetilde T_4,\widetilde T_3\),
satisfy
\begin{equation}\label{pinch:rescaled-errors}
 \|\widetilde T_4\|\le C'\frac{x}{q}=O(x^{-1}),
 \qquad
 \|\widetilde T_3\|\le C'q^{-1/4}=O(x^{-1/2}).
\end{equation}
Indeed, each horizontal slot supplies a factor \(q^{-1/4}\).

The lower bound on \(D\) gives \(E\ge h\).
For an eigenvalue \(d\in[c,C]\) of \(h\), the corresponding ratio is
\[
 \frac Eh=\frac{1+2x}{\sqrt q}
           -\frac{\sqrt q\,d}{\lambda+xd}.
\]
Since
\[
 \frac{1+2x}{\sqrt q}=\sqrt{4+\frac1q}
 \le2+\frac{C_1}{x^2},
 \qquad
 \frac{\sqrt q\,d}{\lambda+xd}
 \ge c_1\frac{x}{\lambda+x},
\]
we have the uniform bound
\begin{equation}\label{pinch:E-bounds}
 h\le E\le
 \left(2+\frac{C_1}{x^2}
          -c_1\frac{x}{\lambda+x}\right)h.
\end{equation}
The constants \(c_1>0,C_1<\infty\) depend only on \(c,C\).

Pass to a subsequence on which \(h,E,\lambda/q\), and the normalized
pair converge.
The limiting forms satisfy \(h\le E\le2h\).
The error contributions tend to zero by
\eqref{pinch:rescaled-errors}.
Lemma~\ref{pinch:algebra} shows that the two limiting summands in
\eqref{pinch:rescaled-leading} are nonnegative.
Since \(F\le0\), both must vanish.
In particular \(P_h(E)=0\) at the limit, so \(a_h,a_V>0\) there and
the mixed blocks of the limiting \(\Xi\) vanish.
The comparison
\[
 c^2a_I\le a_h\le C^2a_I
\]
follows from the corresponding real Gram-matrix bounds.
Hence \(a_I,a_h,a_V\) are bounded below by positive constants along
this subsequence.

We now improve the error estimate on these particular planes.
By the contraction formula \eqref{curv:xi-contraction}, a tensor with exactly three
horizontal indices can pair only one horizontal-horizontal block
with one mixed block.  Thus
\begin{equation}\label{pinch:vanishing-error}
 |\mathcal H_{\widetilde T_3}|
 \le C'q^{-1/4}\|\Xi_{HH}\|\,\|\Xi_{\mathrm{mix}}\|
 =o(x^{-1/2}).
\end{equation}
Here \(\Xi_{HH}\) denotes the block \((\Xi^{i\bar j})\), and
\(\Xi_{\mathrm{mix}}\) comprises the blocks
\((\Xi^{i\bar0})\) and \((\Xi^{0\bar i})\).
The first block is bounded because the normalized pairs form a
compact set, and the mixed blocks tend to zero by the equality case
just proved.
The explicit tensor factor \(q^{-1/4}\) is what makes this an
\(o(x^{-1/2})\) estimate for the evaluated contribution.
No rate of convergence of the planes is being assumed.

Apply \eqref{pinch:square} before taking the limit, with
\[
 \beta=c_1\frac{x}{\lambda+x}-\frac{C_1}{x^2}.
\]
This choice is permitted by \eqref{pinch:E-bounds}, including when
\(\beta<0\).
The positive lower bounds for \(a_I,a_h,a_V\), and boundedness of the
normalized pairs, imply
\begin{equation}\label{pinch:positive-gap}
 L_x\ge c_2\left(\frac{\lambda}{q}
                  +\frac{x}{\lambda+x}\right)-\frac{C_2}{x^2}
\end{equation}
along the subsequence, for some \(c_2>0,C_2<\infty\).
For \(x\ge1\), if \(\lambda\ge x^{3/2}\), then
\[
 \frac{\lambda}{q}\ge\frac{x^{3/2}}{2x^2}
 =\frac1{2\sqrt x}.
\]
If instead \(\lambda\le x^{3/2}\), then
\[
 \frac{x}{\lambda+x}\ge\frac1{\sqrt x+1}
 \ge\frac1{2\sqrt x}.
\]
Thus \eqref{pinch:positive-gap} is bounded below by a positive
constant times \(x^{-1/2}\) for large \(x\).
The evaluated \(\widetilde T_4\) contribution is
\(O(x^{-1})=o(x^{-1/2})\), and
\eqref{pinch:vanishing-error} gives the same little-o comparison for
\(\widetilde T_3\).
It follows that \(F>0\) eventually, a contradiction.

The two regimes exhaust the possible failure sequences.
This proves the asserted finite threshold \(\lambda_0\), uniformly
over all permitted matrices, error tensors, and real planes.
\end{proof}

\subsection{Both curvature bounds at a fixed parameter}

\begin{proof}[Proof of Proposition~\ref{metric:transfer}]
The geometric conclusions concerning the domain and completeness were
proved in Lemma~\ref{metric:complete}.
It remains to obtain both uniform real sectional-curvature bounds.
By Proposition~\ref{curv:tensor}, the errors of the actual metric have
the types and bounds in Lemma~\ref{pinch:positivity}, with
\(C_{\mathrm e}\) depending only on the chart bounds of
Proposition~\ref{metric:transfer}.
Choose any finite \(\lambda\ge\lambda_0\).
Then \(-\mathcal H_R>0\) for every point and every real two-plane.
We show that the sectional-curvature quotient is bounded away from
zero and from negative infinity.

First restrict \(x\) to a bounded interval \([0,X]\).
Normalize the pairs for \(\operatorname{Re}(I+V)\).
The data consisting of \(x\), the matrix \(h\), the component arrays
of \(T_4,T_3\), and the normalized pair range over a closed bounded
subset of a finite-dimensional space: impose
\(cI\le h\le CI\), the stated component types and symmetries, and
the closed norm bounds \(C_{\mathrm e}x,C_{\mathrm e}\sqrt q\).
This set is compact.
The formulas for \(G,D,R\) are continuous on it.
Since Lemma~\ref{pinch:positivity} applies to this entire set,
\(-\mathcal H_R\) has a strictly positive minimum and a finite maximum.
The metric \(G+V\) is uniformly positive definite and bounded there,
so \(a_{G+V}\) is bounded above and away from zero.
Formula~\eqref{curv:sectional} therefore gives both a strictly negative
upper curvature bound and a finite lower curvature bound on
\(0\le x\le X\).

To cover the end \(x\to\infty\), use the horizontal rescaling from
the preceding proof.
Here \(\lambda\) remains fixed: the mixed form \(D/\sqrt q\) will
converge to \(h\), rather than to an arbitrary form between \(h\)
and \(2h\) as allowed in the preceding contradiction argument.
The rescaled Hermitian metric is
\[
 \widetilde g_x=G/\sqrt q+V.
\]
For the fixed finite \(\lambda\),
\begin{equation}\label{pinch:metric-limit}
 \frac G{\sqrt q}-h
 =\frac{\lambda}{\sqrt q}I+
   \left(\frac{x}{\sqrt q}-1\right)h
 =O_\lambda(x^{-1})
\end{equation}
uniformly in the permitted data.
Because \(G\) commutes with \(h\), the formula for \(D\) also gives
\[
 D=xh+\lambda(1+x)hG^{-1}.
\]
Consequently
\begin{equation}\label{pinch:D-limit}
 \frac D{\sqrt q}-h
 =\left(\frac{x}{\sqrt q}-1\right)h+
   \frac{\lambda(1+x)}{\sqrt q}hG^{-1}
 =O_\lambda(x^{-1}).
\end{equation}
The rescaled tensor identity now implies
\begin{equation}\label{pinch:tensor-limit}
 \widetilde R_x=-S_{h+V}+O_\lambda(x^{-1/2})
\end{equation}
in component norm, uniformly over all allowed data.
Indeed, the term \(\lambda S_I\) is divided by \(q\), the term
\(qS_h\) becomes exactly \(S_h\), the mixed term converges by
\eqref{pinch:D-limit}, and the error tensors satisfy
\eqref{pinch:rescaled-errors}.

For the limiting metric \(h+V\) and tensor \(-S_{h+V}\),
\[
 K=-\frac12\left(1+
     3\frac{b_{h+V}^2}{a_{h+V}}\right)
 \in[-2,-1/2].
\]
For an \((I+V)\)-real-orthonormal pair,
\(a_{h+V}\ge\min(c,1)^2\), so these area denominators are uniformly
bounded away from zero.  The forms \(h+V\) also range over a compact
bounded family.  The metric and tensor estimates
\eqref{pinch:metric-limit}--\eqref{pinch:tensor-limit} therefore imply
uniform convergence of the corresponding sectional-curvature
quotients over all real planes.
For sufficiently large \(X_\lambda<\infty\), they give, for example,
\[
 -3\le K_g\le-1/4\qquad(x\ge X_\lambda).
\]
Combining this with the compact-range bounds at \(X=X_\lambda\)
produces finite constants \(0<A\le B\) for which
\[
 -B\le K_g(\sigma)\le-A<0
\]
for every point and every real two-plane.
The tensor identity at \(x=0\), given in
\eqref{curv:zero-section}, is computed in a smooth normalized
fiber frame; equivalently the same bounds extend to the zero section
by smoothness of the metric.
This completes Proposition~\ref{metric:transfer}.
\end{proof}

\bibliographystyle{plain}
\bibliography{references}
\end{document}
